% Root bilingual entry point for the Tufte-style book sample.
% Set the language to chinese and compile this document with XeLaTeX.
\documentclass[nofonts,nohyper,nobib]{tufte-book}

% Allowed values: english, chinese.
\newcommand{\BookLanguage}{chinese}

\usepackage{Thomas-Tufte-bilingual-book}
% 西文字体方案：
%   tufte  = A_Tufte_Style_Book_with_VDQI 字体方案 (Palatino + Helvetica + Bera Mono)。
%   beamer = OUC-Haide-Beamer-Template 字体方案 (RobotoSerif + Menlo + national-2)，
%            从 ./fonts/ 加载（OUC-Haide-Beamer-Template-main/fonts 的副本）。
% 切换只影响西文/等宽/Display 字体，CJK 字体链保持 Songti SC 等不变。
\SetBookFontScheme{tufte}
% \SetBookFontScheme{beamer}
\SetBookLanguage{\BookLanguage}
\OUCSetPalette{0} % 0..23，已扩展为当前对话中讨论过的全部主题配色。
% 如需京都方案中的朱红点睛色，可使用 \textcolor{spotcolor}{...}。

%%
% Book metadata
\title{Thomas Shang 的 Tufte 风格书籍模板}
\date{版本说明}
\author[Thomas Shang \& Tufte-LaTeX Developers]{Thomas Shang \& Tufte-LaTeX Developers}
\publisher{本书出版社}

\begin{document}

% Front matter
\frontmatter

% r.1 blank page
% \blankpage

% v.2 epigraphs
% \newpage\thispagestyle{empty}
% \openepigraph{%
% 公众对糟糕设计往往比对优秀设计更加熟悉。
% 在某种意义上，公众已经被训练去偏爱糟糕设计，
% 因为这正是他们日常所接触的东西。
% 新事物会让人感到威胁，旧事物则让人安心。
% }{Paul Rand}
% \vfill
% \openepigraph{%
% 一个设计者达到完美，并不是因为已经没有东西可以再添加，
% 而是因为已经没有东西可以再删减。
% }{Antoine de Saint-Exup\'{e}ry}
% \vfill
% \openepigraph{%
% \ldots 一个新系统的设计者不仅必须是实现者，
% 也必须是第一个大规模使用者；设计者还应当撰写第一版用户手册\ldots
% 如果我没有亲自参与所有这些工作，
% 数以百计的改进就根本不会发生，
% 因为我既不会想到它们，也不会意识到它们为什么重要。
% }{Donald E. Knuth}


% r.3 full title page
\maketitle


% v.4 copyright page
\newpage
\begin{fullwidth}
~\vfill
\thispagestyle{empty}
\setlength{\parindent}{0pt}
\setlength{\parskip}{\baselineskip}
版权所有 \copyright\ \the\year\ \thanklessauthor

\par 由 \thanklesspublisher 出版

\par tufte-latex.googlecode.com

\par 本文件依据 Apache License, Version 2.0（以下简称“许可协议”）发布；
未经许可，不得使用本文件。你可以在
\url{http://www.apache.org/licenses/LICENSE-2.0}
获得许可协议全文。除非适用法律要求或书面同意，
依据该许可协议分发的软件均按“原样”提供，不附带任何明示或暗示的保证或条件。
有关许可与限制的完整说明，请参见许可协议正文。\index{license}

\par\textit{首次印刷，\monthyear}

\par\medskip
本模板由 Thomas Shang 在 Tufte 风格书籍模板的基础上精心打磨与重新设计而成，
在排版细节、配色体系与中英双语工作流上均做了大量调整与扩展。项目主页：
\url{https://github.com/Hatsuyuki017}。
\end{fullwidth}

% r.5 contents
\tableofcontents

% 当前配色展示页（不编号、不计入目录）。
% 六条 1cm 高的 pal1..pal6 色带，中间间隔 1.6cm。
% 约定（与配色无关）：pal1..pal3 采用黑色字，pal4..pal6 采用白色字。
% 该配色的色彩哲学以 \marginnote 的形式放入边栏。
\cleardoublepage
\thispagestyle{empty}
\chapter*{\OUCPaletteName}
\marginnote{这是模板出厂时的底色。蓝是青年学者的颜色——从清晨第一页讲义、午后图书馆窗台，到深夜实验台的冷光，蓝构成学术工作的整个色温。六层由浅入深的蓝刻意维持理性的克制，模拟一本书从翻开到合上的全过程；它不是一组要被注视的色，而是一组让文字自身被看见的色。}
\noindent\colorbox{oucpal1}{\parbox[c][1cm][c]{\dimexpr\linewidth-2\fboxsep\relax}{\centering\bfseries\color{black}纸海蓝\,Paper Foam\quad\textnormal{\texttt{rgb(239, 246, 255)}}}}\par
\vspace{1.6cm}
\noindent\colorbox{oucpal2}{\parbox[c][1cm][c]{\dimexpr\linewidth-2\fboxsep\relax}{\centering\bfseries\color{black}冰晶蓝\,Desert Foam\quad\textnormal{\texttt{rgb(219, 234, 254)}}}}\par
\vspace{1.6cm}
\noindent\colorbox{oucpal3}{\parbox[c][1cm][c]{\dimexpr\linewidth-2\fboxsep\relax}{\centering\bfseries\color{black}天玻璃蓝\,Sea Smoke\quad\textnormal{\texttt{rgb(96, 165, 250)}}}}\par
\vspace{1.6cm}
\noindent\colorbox{oucpal4}{\parbox[c][1cm][c]{\dimexpr\linewidth-2\fboxsep\relax}{\centering\bfseries\color{white}讲堂蓝\,Crisp Sky\quad\textnormal{\texttt{rgb(37, 99, 235)}}}}\par
\vspace{1.6cm}
\noindent\colorbox{oucpal5}{\parbox[c][1cm][c]{\dimexpr\linewidth-2\fboxsep\relax}{\centering\bfseries\color{white}深思蓝\,Twilight Blue\quad\textnormal{\texttt{rgb(30, 58, 138)}}}}\par
\vspace{1.6cm}
\noindent\colorbox{oucpal6}{\parbox[c][1cm][c]{\dimexpr\linewidth-2\fboxsep\relax}{\centering\bfseries\color{white}铅墨石\,Midnight Hull\quad\textnormal{\texttt{rgb(30, 41, 59)}}}}\par


\listoffigures

\listoftables

% r.7 dedication
\cleardoublepage
~\vfill
\begin{doublespace}
\noindent\fontsize{18}{22}\selectfont\itshape
\nohyphenation
谨献给热爱 \LaTeX{} 的读者，
以及 \mbox{Edward R.~Tufte} 
和 \mbox{Donald E.~Knuth} 的工作欣赏者。
\end{doublespace}
\vfill
\vfill


% r.9 introduction
\cleardoublepage
\chapter*{引言}

本示例书讨论 Thomas Shang 在 Edward Tufte 的几本著作\cite{Tufte2001,Tufte1990,Tufte1997,Tufte2006}
基础上打磨而成的 Tufte 风格书籍模板的设计，以及
\doccls{tufte-book} 与 \doccls{tufte-handout} 文档类的使用方式。

\cleardoublepage

\chapter*{配色索引}

本示例书现在与 beamer 模板共享同一套配色接口，并额外收纳了本项目讨论过的主题配色组。你可以在导言区通过
\doccmd{OUCSetPalette}\{N\} 选择编号配色。每组条目中的六个色块按
\texttt{pal1} 到 \texttt{pal6} 排列，顺序为由浅到深；其中
\texttt{primary}=\texttt{pal5}，\texttt{accent}=\texttt{pal6}。
京都主题另外提供 \texttt{spotcolor}，用于克制使用的朱红点睛色。

\begin{fullwidth}
\scriptsize
\begin{multicols}{3}
\raggedcolumns
\bookpaletteguideentries
\end{multicols}
\end{fullwidth}


%%
% Start the main matter (normal chapters)

\mainmatter

\part{模板参考}

\chapter{环境与盒子演示}
\label{ch:showcase}

本章对照 \texttt{OUC-Haide-Beamer-Template} 的演示帧，让书籍模板与幻灯片模板共享同一套环境名、配色逻辑与代码块样式。每节左侧为环境的\emph{渲染效果}，右侧为对应的\emph{源代码}；当前激活的配色为 \texttt{\OUCPaletteName}，可通过 \doccmd{OUCSetPalette}\{N\} 切换。

\section{配色槽位语义}\index{palette!slots}
六个槽位驱动所有彩色元素：
\texttt{pal1}=\hlred{paleblue} 为代码块底色，
\texttt{pal2} 为 \docenv{Exa}/\docenv{Rmk}/\docenv{Intro} 的正文底色，
\texttt{pal4} 为同三者的标题填充色，
\texttt{pal5}=\hlred{primary} 为强调色（定理标题、章节内强调），
\texttt{pal6}=\hlred{accent} 为最深色（定义标题徽、章标题色）。

\section{中文字符覆盖测试}\index{Chinese character coverage}

下面一段刻意混入若干生僻字与异体字，用于在编译时验证 CJK 字体链能否完整覆盖：\textbf{齉靐龘囙囧囍嬲嫑烎巭玊夯邓阴阳彳亍卩阝廴辶瓦丬灬罒皿夊}。若上述任意字符显示为方框或缺字符号，请检查 \texttt{Songti SC} / \texttt{Source Han Serif SC} / \texttt{Noto Serif CJK SC} 至少有一项可用。局部楷体：\textkai{这是楷体示例，引用古文「天行健，君子以自强不息」。}

\section{定义与定理}\index{environments!De@\texttt{De}}\index{environments!The@\texttt{The}}

\begin{fullwidth}
\begin{minipage}[t]{0.48\linewidth}
\begin{De}{群}{de:g}
群是一个集合，配备一个满足结合律的二元运算，存在单位元，且每个元素都有逆元。
\end{De}
\begin{The}{Lagrange}{the:lag}
若 $H\le G$ 是有限群 $G$ 的子群，则 $|H|$ 整除 $|G|$。
\end{The}
\end{minipage}\hfill
\begin{minipage}[t]{0.48\linewidth}
\begin{lstlisting}[language={[LaTeX]TeX}]
\begin{De}{群}{de:g}
  群是一个集合，配备
  满足结合律的二元运算。
\end{De}
\begin{The}{Lagrange}{the:lag}
  若 $H\le G$，
  则 $|H|$ 整除 $|G|$。
\end{The}
\end{lstlisting}
\end{minipage}
\end{fullwidth}

\section{例与注}\index{environments!Exa@\texttt{Exa}}\index{environments!Rmk@\texttt{Rmk}}

\begin{fullwidth}
\begin{minipage}[t]{0.48\linewidth}
\begin{Exa}{}{exa:cyc}
$\mathbb{Z}/6\mathbb{Z}$ 是阶为 $6$ 的循环群，由 $1$ 的剩余类生成。
\end{Exa}
\begin{Rmk}{}{rmk:tr}
平凡群 $\{e\}$ 是同构意义下阶为 $1$ 的唯一群。
\end{Rmk}
\end{minipage}\hfill
\begin{minipage}[t]{0.48\linewidth}
\begin{lstlisting}[language={[LaTeX]TeX}]
\begin{Exa}{}{exa:cyc}
  $\mathbb{Z}/6\mathbb{Z}$
  是阶为 $6$ 的循环群。
\end{Exa}
\begin{Rmk}{}{rmk:tr}
  平凡群 $\{e\}$ 是
  阶为 $1$ 的唯一群。
\end{Rmk}
\end{lstlisting}
\end{minipage}
\end{fullwidth}

\section{命题、引理与推论}\index{environments!Pro@\texttt{Pro}}\index{environments!Le@\texttt{Le}}\index{environments!Co@\texttt{Co}}

\begin{fullwidth}
\begin{minipage}[t]{0.48\linewidth}
\begin{Pro}{封闭性}{pro:cl}
若 $H,K\le G$，则 $H\cap K$ 对群运算封闭。
\end{Pro}
\begin{Le}{消去律}{le:can}
若群中 $ab=ac$，则 $b=c$。
\end{Le}
\begin{Co}{阶整除}{co:ord}
$|\langle g\rangle|$ 整除 $|G|$。
\end{Co}
\end{minipage}\hfill
\begin{minipage}[t]{0.48\linewidth}
\begin{lstlisting}[language={[LaTeX]TeX}]
\begin{Pro}{封闭性}{pro:cl}
  $H\cap K$ 对群运算封闭。
\end{Pro}
\begin{Le}{消去律}{le:can}
  若 $ab=ac$，则 $b=c$。
\end{Le}
\begin{Co}{阶整除}{co:ord}
  $|\langle g\rangle|$
  整除 $|G|$。
\end{Co}
\end{lstlisting}
\end{minipage}
\end{fullwidth}

\section{证明与注记}\index{environments!Proof@\texttt{Proof}}\index{environments!Boxed@\texttt{Boxed}}

\begin{fullwidth}
\begin{minipage}[t]{0.48\linewidth}
\begin{Proof}
若 $x\in H\cap K$，则 $x\in H$ 且 $x\in K$，故交集仍是子群。
\end{Proof}
\begin{Boxed}
圆角容器，适合用于段落规模的注记、推导小结或视觉上需要聚合的提示信息。
\end{Boxed}
\end{minipage}\hfill
\begin{minipage}[t]{0.48\linewidth}
\begin{lstlisting}[language={[LaTeX]TeX}]
\begin{Proof}
  若 $x\in H\cap K$，
  则 $x\in H$ 且 $x\in K$。
\end{Proof}
\begin{Boxed}
  圆角容器，适合段落
  规模的注记或小结。
\end{Boxed}
\end{lstlisting}
\end{minipage}
\end{fullwidth}

\section{Intro 与 RightIntro}\index{environments!Intro@\texttt{Intro}}\index{environments!RightIntro@\texttt{RightIntro}}

\begin{fullwidth}
\begin{minipage}[t]{0.48\linewidth}
\begin{Intro}{章节引子}{标题徽贴在左上角的引子盒，适合章首或节首预告本节目标。}
\end{Intro}
\begin{RightIntro}{右侧引子}{镜像版本，标题徽贴在右上角。}
\end{RightIntro}
\end{minipage}\hfill
\begin{minipage}[t]{0.48\linewidth}
\begin{lstlisting}[language={[LaTeX]TeX}]
\begin{Intro}{章节引子}
  {标题贴在左上角。}
\end{Intro}
\begin{RightIntro}{右侧引子}
  {标题贴在右上角。}
\end{RightIntro}
\end{lstlisting}
\end{minipage}
\end{fullwidth}

\section{UnderlineBox 与 Quote}\index{UnderlineBox command@\protect\hangleft{\texttt{\tuftebs}}\texttt{UnderlineBox}}\index{Quote command@\protect\hangleft{\texttt{\tuftebs}}\texttt{Quote}}

\begin{fullwidth}
\begin{minipage}[t]{0.48\linewidth}
\UnderlineBox{由上下两条 primary 色短线包裹的下划线块。}
\Quote{数学是科学的女王。}{C.\,F.~高斯}
\end{minipage}\hfill
\begin{minipage}[t]{0.48\linewidth}
\begin{lstlisting}[language={[LaTeX]TeX}]
\UnderlineBox{由两条彩色短线
  包裹的下划线块。}
\Quote{数学是科学的女王。}
      {C.\,F.~高斯}
\end{lstlisting}
\end{minipage}
\end{fullwidth}

\section{Minipage、Eq、EqL}\index{Minipage command@\protect\hangleft{\texttt{\tuftebs}}\texttt{Minipage}}\index{Eq command@\protect\hangleft{\texttt{\tuftebs}}\texttt{Eq}}\index{EqL command@\protect\hangleft{\texttt{\tuftebs}}\texttt{EqL}}

\begin{fullwidth}
\begin{minipage}[t]{0.48\linewidth}
\Minipage{0.46}{\textbf{左侧。}\\双栏布局辅助宏。}{0.46}{\textbf{右侧。}\\两个宽度之和需小于 $1$。}
\Eq{a^2 + b^2 = c^2}
\EqL{e^{i\pi} + 1 = 0}{eq:euler}
\end{minipage}\hfill
\begin{minipage}[t]{0.48\linewidth}
\begin{lstlisting}[language={[LaTeX]TeX}]
\Minipage{0.46}{左侧内容}
         {0.46}{右侧内容}
\Eq{a^2 + b^2 = c^2}
\EqL{e^{i\pi} + 1 = 0}{eq:euler}
\end{lstlisting}
\end{minipage}
\end{fullwidth}

\section{Info / Warn / Result / Goal 盒子}\index{environments!infobox@\texttt{infobox}}\index{environments!warnbox@\texttt{warnbox}}\index{environments!resultbox@\texttt{resultbox}}\index{environments!goalbox@\texttt{goalbox}}

\begin{fullwidth}
\begin{minipage}[t]{0.48\linewidth}
\begin{infobox}
背景、定义与上下文。
\end{infobox}
\begin{warnbox}
注意事项、假设与风险。
\end{warnbox}
\begin{resultbox}
结论、发现与关键产出。
\end{resultbox}
\begin{goalbox}
\centering 简短的高亮信息。
\end{goalbox}
\end{minipage}\hfill
\begin{minipage}[t]{0.48\linewidth}
\begin{lstlisting}[language={[LaTeX]TeX}]
\begin{infobox}...\end{infobox}
\begin{warnbox}...\end{warnbox}
\begin{resultbox}...\end{resultbox}
\begin{goalbox}...\end{goalbox}
\end{lstlisting}
\end{minipage}
\end{fullwidth}

\section{代码块样式}\index{lstlisting style@\texttt{oucbooklisting} style}

\texttt{oucbooklisting} 已注册为 \docpkg{listings} 包的默认样式，并与 beamer 模板共用：paleblue 底色、primary 关键字、accent 字符串、灰色描边。

\subsection{Matlab}
\begin{lstlisting}[language=Matlab]
function y = laplacian(u, dx)
% 均匀网格上的五点离散拉普拉斯算子。
y = (circshift(u,[0 -1]) + circshift(u,[0 1]) ...
   + circshift(u,[-1 0]) + circshift(u,[ 1 0]) - 4*u) / dx^2;
end
\end{lstlisting}

\subsection{Python}
\begin{lstlisting}[language=Python]
def softmax(x):
    """Numerically stable softmax."""
    z = x - x.max(axis=-1, keepdims=True)
    e = np.exp(z)
    return e / e.sum(axis=-1, keepdims=True)
\end{lstlisting}

\subsection{Shell}
\begin{lstlisting}[language=bash]
latexmk -xelatex -interaction=nonstopmode Thomas-Tufte-book-zh.tex
makeindex Thomas-Tufte-book-zh.idx
latexmk -xelatex -interaction=nonstopmode Thomas-Tufte-book-zh.tex
\end{lstlisting}

\subsection{LaTeX}
\begin{lstlisting}[language={[LaTeX]TeX}]
\usepackage{Thomas-Tufte-bilingual-book}
\SetBookLanguage{chinese}
\OUCSetPalette{0}
\begin{document}
  \maketitle
  \tableofcontents
  \mainmatter
\end{document}
\end{lstlisting}

\cleardoublepage


\part{第一部分}

\chapter{Tufte 书籍的设计}
\label{ch:tufte-design}


\newthought{一本书的页面}通常会被划分为三个主要部分：前言部分（front matter，也称 preliminary matter 或 prelim）、正文部分（main matter，即核心正文）和后记部分（back matter，也称 end matter）。

\newthought{前言部分}是指主文之前出现的所有材料。下表列出了 \VDQI、\EI、\VE 和 \BE 四本书在前言中出现的内容及其页码。括号中的页码表示该页没有印刷页码，但仍然被计入页面顺序之中。

\bigskip
\begin{minipage}{\textwidth}
\begin{center}
\begin{tabular}{lcccc}
\toprule
 & \multicolumn{4}{c}{书籍} \\
\cmidrule(l){2-5}
\booktableheadrow
\booktableheadcell{页面内容} & \booktableheadcell{\vdqi} & \booktableheadcell{\ei} & \booktableheadcell{\ve} & \booktableheadcell{\be} \\
\midrule
空白半标题页 & \hangp{1} & \hangp{1} & \hangp{1} & \hangp{1} \\
卷首插页\footnotemark{}
  & \hangp{2} & \hangp{2} & \hangp{2} & \hangp{2} \\
完整标题页 & \hangp{3} & \hangp{3} & \hangp{3} & \hangp{3} \\
版权页 & \hangp{4} & \hangp{4} & \hangp{4} & \hangp{4} \\
目录 & \hangp{5} & \hangp{5} & \hangp{5} & \hangp{5} \\
%Blank & -- & \hangp{6} & \hangp{6} & \hangp{6} \\
题献页 & \hangp{6} & \hangp{7} & \hangp{7} & 7 \\
%Blank & -- & \hangp{8} & -- & \hangp{8} \\
题辞页 & -- & -- & \hangp{8} & -- \\
引言 & \hangp{7} & \hangp{9} & \hangp{9} & 9 \\
\bottomrule
\end{tabular}
\end{center}
\end{minipage}
\vspace{-7\baselineskip}\footnotetext{这一页的内容因书而异。在
  \vdqi 中该页为空白；在 \ei 和 \ve 中该页为卷首插页；
  在 \be 中则包含三则题辞。}
\vspace{7\baselineskip}

\bigskip
Tufte 书籍中的前言设计与传统前言设计略有不同。首先，传统书籍的前言页码通常采用小写罗马数字（例如 \eg, i, ii, iii, iv,~\ldots）。其次，前言部分的页码序列通常与正文页码分开，也就是说正文开始时页码会重新从 1 计起。相对地，Tufte 使用的是阿拉伯数字，而且其页码计数与正文保持在同一序列之中。

Tufte 的四本书在设计上也存在一些差异。标题页对开的那一页，也就是上表中的 “frontispiece”，在不同书中内容不同。在 \VDQI 中这一页为空白；在 \EI 和 \VE 中这一页放置卷首插图；而在 \BE 中这一页则包含三则题辞。

在 \vdqi 中，题献页位于第 6 页，与引言相对；而在其他几本书中，题献页单独占据一个跨页。在 \ve 中，一则题辞与引言首页共处于同一跨页中。

前言部分的页码通常都不印出，只有在 \be 中，页码从题献页开始出现。

\newthought{完整标题页}在 Tufte 的几本书之间也略有差别。四本书都把作者名放在页首，把书名置于页面中线稍上方，把出版社放在下边距附近。它们之间的一些差异归纳如下表。

\bigskip
\begin{center}
\footnotesize
\begin{tabular}{lllll}
\toprule
\booktableheadrow
\booktableheadcell{特征} & \booktableheadcell{\vdqi} & \booktableheadcell{\ei} & \booktableheadcell{\ve} & \booktableheadcell{\be} \\
\midrule
作者 & & & & \\
\quad 字体类别 & serif   & serif   & serif   & sans serif \\
\quad 样式    & italics & italics & italics & upright, caps \\
\quad 字号     & 24 pt   & 20 pt   & 20 pt   & 20 pt \\
\addlinespace
标题 & & & & \\
\quad 字体类别 & serif   & serif   & serif   & sans serif \\
\quad 样式    & upright & italics & upright & upright, caps \\
\quad 字号     & 36 pt   & 48 pt   & 48 pt   & 36 pt \\
\addlinespace
副标题 & & & & \\
\quad 字体类别 & \na     & \na     & serif   & \na \\
\quad 样式    & \na     & \na     & upright & \na \\
\quad 字号     & \na     & \na     & 20 pt   & \na \\
\addlinespace
版本说明 & & & & \\
\quad 字体类别 & sans serif    & \na  & \na  & \na \\
\quad 样式    & upright, caps & \na  & \na  & \na \\
\quad 字号     & 14 pt         & \na  & \na  & \na \\
\addlinespace
出版社 & & & & \\
\quad 字体类别 & serif   & serif   & serif   & sans serif \\
\quad 样式    & italics & italics & italics & upright, caps \\
\quad 字号     & 14 pt   & 14 pt   & 14 pt   & 14 pt \\
\bottomrule
\end{tabular}
\end{center}

\begin{figure*}[p]
\fbox{\includegraphics[width=0.45\linewidth]{vdqi-title.pdf}}
\hfill
\fbox{\includegraphics[width=0.45\linewidth]{ei-title.pdf}}
\\\vspace{\baselineskip}
\fbox{\includegraphics[width=0.45\linewidth]{ve-title.pdf}}
\hfill
\fbox{
\includegraphics[width=0.45\linewidth]{be-title.pdf}}
\end{figure*}

\newthought{目录页}让我们第一次真正看到正文部分的结构。\VDQI 被分成两个部分，每个部分包含若干章节；而另外三本书只有章节结构，并不分部。

\begin{figure*}[p]\index{table of contents}
\fbox{\includegraphics[width=0.45\linewidth]{vdqi-contents.pdf}}
\hfill
\fbox{\includegraphics[width=0.45\linewidth]{ei-contents.pdf}}
\\\vspace{\baselineskip}
\fbox{\includegraphics[width=0.45\linewidth]{ve-contents.pdf}}
\hfill
\fbox{\includegraphics[width=0.45\linewidth]{be-contents.pdf}}
\end{figure*}


\section{字体}\label{sec:typefaces1}\index{typefaces}
\index{fonts|see{typefaces}}

Tufte 的书主要使用两种字体：Bembo 和 Gill Sans。Bembo 用于标题和正文，而 Gill Sans 则用于 \BE 的标题页和开篇题辞。

由于默认的 \LaTeX{} 安装中通常没有 Bembo 和 Gill Sans，\TL 文档类默认分别使用 Palatino 和 Helvetica 作为替代。此外，Bera Mono 用于 \texttt{等宽文本} 的排版。

\TL 文档类中定义的字号如下：

\noindent 正文默认字号为 10/14，边注与题注为 8/10；主要标题依次上升到
11/15 与 12/16，而最大的书籍标题在章节处达到 20/30，在部分标题处达到
24/36。

\section{标题层级}\label{sec:headings1}\index{headings}

Tufte 的书中包含以下标题层级：部分、章节，\sidenote{部分和章节只在 \texttt{tufte-book} 类中定义。} 以及 sections、subsections 和 paragraphs。默认未定义的层级包括 sub-subsections 和 subparagraphs。

\begin{table}[h]
  \begin{center}
    \footnotesize%
    \begin{tabular}{lcr}
      \toprule
      \booktableheadrow
      \booktableheadcell{标题层级} & \booktableheadcell{样式} & \booktableheadcell{尺寸} \\
      \midrule
      Part & roman & \measure{24}{36}{40} \\
      Chapter & italic & \measure{20}{30}{40} \\
      Section & italic & \measure{12}{16}{26} \\
      Subsection & italic & \measure{11}{15}{26} \\
      Paragraph & italic & 10/14 \\
      \bottomrule
    \end{tabular}
  \end{center}
  \caption{\BE 中使用的标题样式。}
  \label{tab:heading-styles}
\end{table}

\paragraph{Paragraph} 段落标题像这里这样使用斜体引导，并与随后的正文之间保留少量空白。

\section{环境}

以下特征定义了不同环境的排版方式：


\begin{table}[h]
  \begin{center}
    \footnotesize%
    \begin{tabular}{lcl}
      \toprule
      \booktableheadrow
      \booktableheadcell{环境} & \booktableheadcell{字号} & \booktableheadcell{说明} \\
      \midrule
      正文 & \measure{10}{14}{26} & \\
      块引用 & \measure{9}{12}{24} & 左右各缩进 \unit[1]{pc} \\
      边注 & \measure{8}{10}{12} & 边注编号内联排版，后接一个词距 \\
      题注 & \measure{8}{10}{12} &  \\
      \bottomrule
    \end{tabular}
  \end{center}
  \caption{\BE 中使用的环境样式。}
  \label{tab:environment-styles}
\end{table}



\chapter[关于 tufte-book 文档类的使用]{关于 \texttt{tufte-book} 文档类的使用}
\label{ch:tufte-book}

\TL 文档类定义了一种与 Edward Tufte 在其书籍和讲义中所使用风格相近的排版样式。这种风格以大量使用边注、图文紧密结合以及讲究的字体排版而著称。本文档既是 \TL 文档类功能的展示，也是其使用方式的风格指南。

\section{页面布局}\label{sec:page-layout}
\subsection{标题}\label{sec:headings}\index{headings}
这种风格提供 \textsc{a}- 和 \textsc{b}-heads，也就是上文示范过的 \Verb|\section| 和 \Verb|\subsection|。

如果你需要超过两个层级的节标题，目前需要自己定义更低层级的样式；因为在默认设置中，低于 \Verb|\subsection| 的标题样式并未预先定义。正如 Bringhurst 在 \textit{The Elements of Typographic Style} 一书中指出的那样，\cite{Bringhurst2005} 你应当“使用你所需要的标题层级，不多也不少”。

如果尝试使用 \Verb|\subsubsection| 或更低层级标题，\TL 类会报错。

% let's start a new thought -- a new section
\newthought{在他后期的书中}，\cite{Tufte2006} Tufte 在每个 section 的开头加入一点竖直空白，取消首段缩进，并将句子的前几个词设置为 \textsc{small caps}。若要在本样式中实现这种效果，请使用 \doccmddef{newthought} 命令：
\begin{docspec}
  \doccmd{newthought}\{在他后期的书中\}, Tufte starts\ldots
\end{docspec}


\section{边注}\label{sec:sidenotes}
这种风格最显著、最具辨识度的特征之一就是对边注的大量使用。页面预留了较宽的边栏，以容纳边注和小型插图。任何 \doccmd{footnote} 都会自动转换成边注。\footnote{这是通过 \texttt{\textbackslash footnote} 命令输入的一个边注。} 如果你想在边栏中放置信息，但不需要边注标记，也就是正文中的上标数字，可以使用 \doccmd{marginnote} 命令。\marginnote{这是一个边栏注。注意它前面没有编号，正文中对应位置也没有编号提示。}

\doccmddef{sidenote} 命令的规范如下：
\begin{docspec}
  \doccmd{sidenote}[\docopt{number}][\docopt{offset}]\{\docarg{Sidenote text.}\}
\end{docspec}

\docopt{number} 和 \docopt{offset} 两个参数都是可选的。如果提供了 \docopt{number}，该数字会被用作当前边注的编号，但只会影响当前这一条边注，不会改变后续边注的编号序列。

有时边注可能会压到边栏中的其他文字或图形。这种情况下，你可以通过 \docopt{offset} 参数提供一个长度值，来调整边注的垂直位置。有效的长度值例如：
\begin{docspec}
  \ttfamily 1.0in \qquad 2.54cm \qquad 254mm \qquad 6\Verb|\baselineskip|
\end{docspec}
如果该值为正，则边注会向下移动；如果为负，则边注会上移。

虽然 \docopt{number} 和 \docopt{offset} 都是可选参数，但必须按顺序提供。如果你只想调整边注的垂直位置而保持默认编号，可以写成：
\begin{docspec}
  \doccmd{sidenote}[][\docopt{offset}]\{\docarg{Sidenote text.}\}
\end{docspec}
空的第一组方括号告诉 \Verb|\sidenote| 命令继续使用默认边注编号。

如果你 \emph{只} 想改变边注编号，则可以完全省略 \docopt{offset}：
\begin{docspec}
  \doccmd{sidenote}[\docopt{number}]\{\docarg{Sidenote text.}\}
\end{docspec}

\doccmddef{marginnote} 命令也接受一个类似的 \docarg{offset} 参数：
\begin{docspec}
  \doccmd{marginnote}[\docopt{offset}]\{\docarg{Margin note text.}\}
\end{docspec}

\section{参考文献}
参考文献也会与其引用并排放入边注中。这可以通过普通的 \doccmddef{cite} 命令实现。\sidenote{本文第一段中已经包含了一个引用示例。}

完整参考文献列表也可以通过 \doccmddef{bibliography} 命令自动打印。（本文末尾会给出示例。）如果你不希望在文末打印参考文献表，则可以使用 \doccmddef{nobibliography} 命令代替。

若要在同一位置插入多个引用，\cite[-3\baselineskip]{Tufte2006,Tufte1990} 可以提供一个由逗号分隔的键列表，并同样使用可选的垂直偏移参数：\Verb|\cite{Tufte2006,Tufte1990}|。
\begin{docspec}
  \doccmd{cite}[\docopt{offset}]\{\docarg{bibkey1,bibkey2,\ldots}\}
\end{docspec}

\section{图与表}\label{sec:figures-and-tables}
图像和图形在 Tufte 的作品中发挥着核心作用。除了标准的 \docenvdef{figure} 和 \docenvdef{tabular} 环境之外，这一风格还提供了适用于通栏浮动体的特殊 figure 和 table 环境。

整页宽度的图表可以放在 \docenvdef{figure*} 或 \docenvdef{table*} 环境中。若要把图或表放到边栏中，可以使用 \docenvdef{marginfigure} 或 \docenvdef{margintable} 环境，如下所示，见图~\ref{fig:marginfig}：

\begin{marginfigure}%
  \includegraphics[width=\linewidth]{helix}
  \caption{这是一个边栏图。该螺旋线由
    $x = \cos(2\pi z)$、$y = \sin(2\pi z)$ 和 $z = [0, 2.7]$ 定义。
    该图使用 Asymptote 绘制，见 \protect\url{http://asymptote.sf.net/}。}
  \label{fig:marginfig}
\end{marginfigure}

\begin{docspec}
\textbackslash begin\{marginfigure\}\\
  \qquad\textbackslash includegraphics\{helix\}\\
  \qquad\textbackslash caption\{This is a margin figure.\}\\
  \qquad\textbackslash label\{fig:marginfig\}\\
\textbackslash end\{marginfigure\}\\
\end{docspec}

\docenv{marginfigure} 和 \docenv{margintable} 环境都接受一个可选参数 \docopt{offset}，用于调整图或表的垂直位置。相关示例可参考上文 “\nameref{sec:sidenotes}” 一节。其规范如下：
\begin{docspec}
  \textbackslash{begin\{marginfigure\}[\docopt{offset}]}\\
  \qquad\ldots\\
  \textbackslash{end\{marginfigure\}}\\
  \mbox{}\\
  \textbackslash{begin\{margintable\}[\docopt{offset}]}\\
  \qquad\ldots\\
  \textbackslash{end\{margintable\}}\\
\end{docspec}

图~\ref{fig:fullfig} 演示了 \docenv{figure*} 环境，而图~\ref{fig:textfig} 则演示了普通的 \docenv{figure} 环境。

\begin{figure*}[h]
  \includegraphics[width=\linewidth]{sine.pdf}%
  \caption{这幅图展示了 $y = \sin x$ 在大约 $x = [-10, 10]$ 区间上的图像。
  \emph{请注意，这幅图占据了整页宽度。}}%
  \label{fig:fullfig}%
\end{figure*}

\begin{figure}
  \includegraphics{hilbertcurves.pdf}
%  \checkparity This is an \pageparity\ page.%
  \caption[不同阶数 $n$ 的 Hilbert 曲线。][5\baselineskip]{不同阶数 $n$ 的 Hilbert 曲线。\emph{请注意，这幅图只占据正文版心宽度。}}
  \label{fig:textfig}
  %\zsavepos{pos:textfig}
\end{figure}

与边注和边栏注类似，图表题注有时也需要调整垂直位置。\doccmddef{caption} 命令现在接受第二个可选参数，用于通过给定一个长度值来实现这种调整。你可以采用以下任意一种形式来书写题注：
\begin{docspec}
  \doccmd{caption}\{\docarg{long caption}\}\\
  \doccmd{caption}[\docarg{short caption}]\{\docarg{long caption}\}\\
  \doccmd{caption}[][\docopt{offset}]\{\docarg{long caption}\}\\
  \doccmd{caption}[\docarg{short caption}][\docopt{offset}]%
                  \{\docarg{long caption}\}
\end{docspec}
正的 \docopt{offset} 会把题注向页面下方推移。若提供了 short caption，它将出现在图表目录中；否则目录中会显示长题注。需要注意的是，虽然 \docopt{short caption} 和 \docopt{offset} 都是可选参数，但它们必须按顺序提供。因此，如果你想指定 \docopt{offset} 却不提供 \docopt{short caption}，就必须先写一组空方括号 \Verb|[]|，以通知 \doccmd{caption} 使用默认的长题注。例如，上图~\ref{fig:textfig} 的题注写法是：
\begin{docspec}
  \doccmd{caption}[Hilbert curves...][6pt]\{Hilbert curves...\}
\end{docspec}

表~\ref{tab:normaltab} 展示了使用 \docpkg{booktabs} 宏包排出的表格。请注意其中没有竖线，因为竖线只会让表格数据显得更加杂乱。

\begin{table}[ht]
  \centering
  \fontfamily{ppl}\selectfont
  \begin{tabular}{ll}
    \toprule
    \booktableheadrow
    \booktableheadcell{边距项目} & \booktableheadcell{长度} \\
    \midrule
    纸张宽度 & \unit[8\nicefrac{1}{2}]{inches} \\
    纸张高度 & \unit[11]{inches} \\
    正文版心宽度 & \unit[6\nicefrac{1}{2}]{inches} \\
    正文与边注间距 & \unit[\nicefrac{3}{8}]{inches} \\
    边注宽度 & \unit[2]{inches} \\
    \bottomrule
  \end{tabular}
  \caption{这是 Tufte-handout 类中各类边距的尺寸。}
  \label{tab:normaltab}
  %\zsavepos{pos:normaltab}
\end{table}

\newthought{有时} \LaTeX{} 会给出如下错误信息：\label{err:too-many-floats}
\begin{docspec}
  Error: Too many unprocessed floats
\end{docspec}
\LaTeX{} 会尝试把浮动体放置到页面上最佳的位置。但在页面尚未排版完成之前，它并不知道那些最佳位置究竟在哪里。如果某一页上的浮动体太多，包括边注、边栏注、图、表等，\LaTeX{} 可能会因为用于追踪浮动体的“槽位”不够而报出上述错误。

\LaTeX{} 初始只分配 18 个槽位来存储浮动体。为了解决这个限制，\TL 文档类提供了 \doccmddef{morefloats} 命令，用于额外保留更多槽位。

第一次调用 \doccmd{morefloats} 时，会额外分配 34 个槽位；第二次调用时，会再额外分配 26 个槽位。

\doccmd{morefloats} 最多只能调用两次。第三次调用时就会触发错误。这是因为再继续分配更多浮动体槽位可能会使 \LaTeX{} 内存不足。

如果在调用 \doccmd{morefloats} 两次之后，你仍然遇到 \texttt{Too many unprocessed floats} 错误，那么还有几个办法可以尝试。

\doccmddef{FloatBarrier} 命令会立即处理当前所有浮动体，然后再继续排版后续内容。由于 \doccmd{FloatBarrier} 会开启一个新段落，因此你应当把它放在段落的开头或结尾。

\doccmddef{clearpage} 命令也会先处理所有浮动体，但它不是开启新段落，而是直接另起一页。

你也可以稍微调整浮动体的位置，例如把某个图或表移到下一页，或者减少边注数量。每一个边注实际上会占用 \emph{两个} 槽位。

当浮动体被放置完成之后，\LaTeX{} 会把相应槽位标记为空闲，以供下一页继续使用。

\section{题注}
你可能会注意到题注有时会出现对齐不佳的问题。由于 \LaTeX{} 的浮动机制，我们无法确切知道浮动体最终会被放到哪里。因此，\TL 文档类提供了一些命令来手动覆盖题注的位置。

\paragraph{垂直对齐} 若要覆盖垂直对齐方式，请在浮动体环境内部使用 \doccmd{setfloatalignment} 命令。例如：

\begin{fullwidth}
\begin{docspec}
  \textbackslash begin\{figure\}[btp]\\
  \qquad \textbackslash includegraphics\{sinewave\}\\
  \qquad \textbackslash caption\{This is an example of a sine wave.\}\\
  \qquad \textbackslash label\{fig:sinewave\}\\
  \qquad \hlred{\textbackslash setfloatalignment\{b\}\% forces caption to be bottom-aligned}\\
  \textbackslash end\{figure\}
\end{docspec}
\end{fullwidth}

\noindent \doccmddef{setfloatalignment} 命令的语法如下：

\begin{docspec}
  \doccmd{setfloatalignment}\{\docopt{pos}\}
\end{docspec}

\noindent 其中 \docopt{pos} 可以取 \texttt{b}，即底部对齐题注，或 \texttt{t}，即顶部对齐题注。

\paragraph{水平对齐}\label{par:overriding-horizontal}
若要覆盖水平对齐方式，可以在浮动体环境中使用 \doccmd{forceversofloat} 或 \doccmd{forcerectofloat} 命令。例如：

\begin{fullwidth}
\begin{docspec}
  \textbackslash begin\{figure\}[btp]\\
  \qquad \textbackslash includegraphics\{sinewave\}\\
  \qquad \textbackslash caption\{This is an example of a sine wave.\}\\
  \qquad \textbackslash label\{fig:sinewave\}\\
  \qquad \hlred{\textbackslash forceversofloat\% forces caption to be set to the left of the float}\\
  \textbackslash end\{figure\}
\end{docspec}
\end{fullwidth}

\doccmddef{forceversofloat} 命令会让算法假定该浮动体被放在左页，也就是 verso 页；相反，\doccmddef{forcerectofloat} 会让算法假定该浮动体位于右页，也就是 recto 页。


\section{通栏文本块}

除了新的浮动体类型之外，这一风格还提供了一个 \docenvdef{fullwidth} 环境，可让内容横跨正文版心与边注区域。

\begin{Verbatim}
\begin{fullwidth}
Lorem ipsum dolor sit amet...
\end{fullwidth}
\end{Verbatim}

\begin{fullwidth}
\small\itshape\lipsum[1]
\end{fullwidth}

\section{排版学}\label{sec:typography}

\subsection{字体}\label{sec:typefaces}\index{typefaces}
如果系统中安装了 Palatino、\textsf{Helvetica} 和 \texttt{Bera Mono}，本样式会自动使用它们。否则将退回使用 Computer Modern 字体族。

\subsection{字距调整}\label{sec:letterspacing}
这个文档类新增了两个字距调整命令，并改进了若干现有命令的字距行为。

在排版 \allcaps{ALL CAPS} 或 \smallcaps{small caps} 这样的字符串时，字母之间的间距应略微增大。\cite[][p.32]{Bringhurst2005} \doccmddef{allcaps} 命令可以为 \allcaps{FULL CAPITAL LETTERS} 提供正确的字距，而 \doccmddef{smallcaps} 命令则会为 \smallcaps{small capital letters} 提供合适的字距。这两个命令也会自动将文本转换成相应的大小写形式。

\doccmddef{textsc} 命令也被重新定义，以包含字距调整。不过，\doccmd{textsc} 参数中的字母大小写会保持原样。因此，你可以同时使用大写和小写字母：
\textsc{The Initial Letters Of The Words In This Sentence Are Capitalized.}



\section{文档类选项}\label{sec:options}

\index{class options|(}
\doccls{tufte-book} 类基于 \LaTeX\ 的 \doccls{book} 文档类，因此你可以传入所有常见的 book 选项。不过，\doccls{tufte-book} 也提供了一些专属选项。

\docclsoptdef{a4paper} 选项会把纸张大小设为 \smallcaps{A4}，而不是默认的 \smallcaps{US} letter。

\docclsoptdef{sfsidenotes} 选项会把边注和标题区排成 \textsf{sans serif} 字体，而不是默认的 roman。

\docclsoptdef{twoside} 选项会修改页眉，使页码显示在外侧边缘，而不是像 \docclsoptdef{oneside} 模式那样总是出现在右侧边缘。

\docclsoptdef{symmetric} 选项会把边注排在页面外侧边缘。这更符合传统书籍的印刷方式，但与 Tufte 自己的书籍设计不同，后者总是把边注排在右侧。该选项会隐式启用 \docclsopt{twoside}。

\docclsoptdef{justified} 选项会把正文设置为左右两端对齐。默认状态是右侧参差。Tufte 的书正文就是这样排的；它能减少不必要的断词，也使较窄的文本栏更易阅读。

\docclsoptdef{bidi} 选项会加载 \docpkg{bidi} 宏包，用于配合 \tXeLaTeX\ 排版双向文本。由于 \docpkg{bidi} 必须在定义 sidenotes 和 cite 命令之前就被加载，因此不能在文档前导区中后置加载。

\docclsoptdef{debug} 选项会让 \TL 类把调试信息输出到日志文件中，便于排查错误，同时还会把图像替换成轮廓框。

\docclsoptdef{nofonts} 选项会阻止 \TL 类自动加载 Palatino 和 Helvetica。如果你希望自行加载字体，应当使用这个选项。如果你正在使用 \tXeLaTeX，实际上该选项是隐含启用的，也就是说，使用 \tXeLaTeX 时 Palatino 和 Helvetica 不会自动加载。

\docclsoptdef{nols} 选项会关闭字距处理代码。\TL 类会尝试加载合适的字距宏包，也就是 pdf\TeX 的 \docpkg{letterspace} 或 \docpkg{soul}。不过如果你是在 \tXeLaTeX\ 下并搭配 \docpkg{fontenc} 使用，就需要自己配置字距。

\docclsoptdef{notitlepage} 选项会让 \doccmd{maketitle} 生成标题块而不是完整标题页。\doccls{book} 类默认使用标题页，\doccls{handout} 类默认使用标题块。与之对应的 \docclsoptdef{titlepage} 选项则会强制生成完整标题页。

\docclsoptdef{notoc} 选项会关闭 \TL 自定义的目录，也就是 \textsc{toc} 设计。目前这套目录设计只显示无编号章节标题，不显示 sections 或 subsections。启用 \docclsopt{notoc} 之后，会退回到 \LaTeX{} 默认目录样式。

\docclsoptdef{nohyper} 选项会阻止加载 \docpkg{hyperref} 宏包。默认情况下会加载 \docpkg{hyperref}，并把 \doccmd{title} 与 \doccmd{author} 的内容写入生成 \textsc{pdf} 的元数据中。

\index{class options|)}




\chapter[定制 Tufte-LaTeX]{定制 \TL}
\label{ch:customizing}

默认情况下，\TL 文档类会尽可能精确地模拟 Tufte 书籍的设计。不过每一份文档都有自己的需求，因此你可能会遇到默认设置不足以胜任的情况。本章将探讨多种调整 \TL 文档类的方法，以便更好地适应你的文档需求。

\section{文件钩子}
\label{sec:filehooks}

\index{file hooks|(}
如果你经常使用 \TL 文档类来制作文档，那么把自定义设置保存在单独文件中，会比每次都复制到前导区更方便。\TL 文档类提供了三个文件钩子：\docfilehook{tufte-common-local.tex}{common}、\docfilehook{tufte-book-local.tex}{book} 和 \docfilehook{tufte-handout-local.tex}{handout}。\sloppy

\begin{description}
  \item[\docfilehook{tufte-common-local.tex}{common}]
    如果该文件存在，它会在所有 \TL 文档类中被加载，并且位置位于任何文档类特定代码之前。如果你的自定义既适用于 book 类，也适用于 handout 类，就应当使用这个公共钩子文件。
  \item[\docfilehook{tufte-book-local.tex}{book}]
    如果该文件存在，它会在所有 common 和 book 特定代码之后加载。如果你的自定义只针对 book 类，请使用这个钩子文件。
  \item[\docfilehook{tufte-common-handout.tex}{handout}]
    如果该文件存在，它会在所有 common 和 handout 特定代码之后加载。如果你的自定义只适用于 handout 类，请使用这个钩子文件。
\end{description}

\index{file hooks|)}

\section{带编号的节标题}
\label{sec:numbered-sections}
\index{headings!numbered}

虽然 Tufte 在自己的书中通常不使用编号标题，但如果你有需求，也可以通过改变 \doccounter{secnumdepth} 计数器的值来启用编号。请根据下表选择你希望停止编号的标题层级，并将 \doccounter{secnumdepth} 设为该值。例如，如果你希望部分和章节带编号，而 section 和 subsection 不带编号，则可以使用：
\begin{docspec}
  \doccmd{setcounter}\{secnumdepth\}\{0\}
\end{docspec}

\TL 文档类默认的 \doccounter{secnumdepth} 值为 $-1$。

\begin{table}
  \footnotesize
  \begin{center}
    \begin{tabular}{lr}
      \toprule
      \booktableheadrow
      \booktableheadcell{标题层级} & \booktableheadcell{值} \\
      \midrule
      Part（在 \doccls{tufte-book} 中） & $-1$ \\
      Part（在 \doccls{tufte-handout} 中） & $0$ \\
      Chapter（仅在 \doccls{tufte-book} 中） & $0$ \\
      Section & $1$ \\
      Subsection & $2$ \\
      Subsubsection & $3$ \\
      Paragraph & $4$ \\
      Subparagraph & $5$ \\
      \bottomrule
    \end{tabular}
  \end{center}
  \caption{与 \texttt{secnumdepth} 计数器配合使用的标题层级。}
\end{table}

\section{改变纸张尺寸}
\label{sec:paper-size}

\TL 类目前只提供三种纸张规格：\textsc{a4}、\textsc{b5} 和 \textsc{us} letter。如果你希望指定不同的纸张尺寸，以及页边距，可以在文档前导区或钩子文件中使用 \doccmd[geometry]{geometrysetup} 命令。该命令的完整说明请参见 \docpkg{geometry} 宏包文档。\cite[-1em][p.33]{pkg-geometry}


\section{定制边栏材料}
\label{sec:marginal-material}

边栏材料包括边注、引用、边栏注和题注。通常情况下，边栏材料的对齐方式会跟随正文。如果你指定了 \docclsopt{justified} 选项，那么所有边栏材料也会采用两端对齐；如果不指定，则边栏材料默认右侧参差。

你也可以通过以下文档类选项单独设置边栏材料的对齐方式：\docclsopt{sidenote}、\docclsopt{marginnote}、\docclsopt{caption}、\docclsopt{citation} 和 \docclsopt{marginals}。每个选项都对应其名称所指的边栏材料类型，而 \docclsopt{marginals} 会同时设置这四类材料的对齐方式。

每个选项都可以取以下五种对齐方式之一：
\begin{description}
  \item[\docclsopt{justified}] 左右两端对齐，也就是同时齐左齐右。
  \item[\docclsopt{raggedleft}] 始终左侧参差，不考虑该材料位于哪一页。
  \item[\docclsopt{raggedright}] 始终右侧参差，不考虑该材料位于哪一页。
  \item[\doccls{raggedouter}] 如果材料位于跨页左页，也就是 verso，则左侧参差；否则右侧参差。这个选项在配合 \docclsopt{symmetric} 使用时尤其有用。
  \item[\docclsopt{auto}] 如果指定了 \docclsopt{justified}，则采用两端对齐；否则采用右侧参差。这也是未显式指定时的默认值。
\end{description}

\noindent 例如，
\begin{docspec}
  \doccmdnoindex{documentclass}[symmetric,justified,marginals=raggedouter]\{tufte-book\}
\end{docspec}
会把正文设为两端对齐，并让所有边栏材料，包含边注、边栏注、题注和引用，紧贴正文一侧，而外侧边缘保持参差。

\newthought{边栏材料的字体与样式}也可以通过以下命令来修改：

\begin{docspec}
  \doccmd{setsidenotefont}\{\docopt{font commands}\}\\
  \doccmd{setcaptionfont}\{\docopt{font commands}\}\\
  \doccmd{setmarginnotefont}\{\docopt{font commands}\}\\
  \doccmd{setcitationfont}\{\docopt{font commands}\}
\end{docspec}

\doccmddef{setsidenotefont} 用于设置边注字体，\doccmddef{setcaptionfont} 用于设置题注字体，\doccmddef{setmarginnotefont} 用于设置边栏注字体，\doccmddef{setcitationfont} 用于设置引用字体。\docopt{font commands} 中可以包含字号命令，例如 \doccmdnoindex{footnotesize}、\doccmdnoindex{Huge}，字体样式命令，例如 \doccmdnoindex{sffamily}、\doccmdnoindex{ttfamily}、\doccmdnoindex{itshape}，颜色命令，例如 \doccmdnoindex{color}\texttt{\{blue\}}，以及其他调整方式。

例如，如果你希望题注使用斜体无衬线字体，可以写成：
\begin{docspec}
  \doccmd{setcaptionfont}\{\doccmdnoindex{itshape}\doccmdnoindex{sffamily}\}
\end{docspec}


\chapter{兼容性问题}
\label{ch:compatibility}

当你把已有文档从其他文档类切换到 \TL 文档类时，往往需要做一些调整。

\section{从 \doccls{article} 转换到 \doccls{tufte-handout}}

以下 \doccls{article} 选项不受支持：\docclsopt{10pt}、\docclsopt{11pt}、\docclsopt{12pt}、\docclsopt{a5paper}、\docclsopt{b5paper}、\docclsopt{executivepaper}、\docclsopt{legalpaper}、\docclsopt{landscape}、\docclsopt{onecolumn} 和 \doccls{twocolumn}。

以下标题命令不受支持：\doccmd{subsubsection} 和 \doccmd{subparagraph}。

\section{从 \doccls{book} 转换到 \doccls{tufte-book}}

以下 \doccls{report} 选项不受支持：\docclsopt{10pt}、\docclsopt{11pt}、\docclsopt{12pt}、\docclsopt{a5paper}、\docclsopt{b5paper}、\docclsopt{executivepaper}、\docclsopt{legalpaper}、\docclsopt{landscape}、\docclsopt{onecolumn} 和 \doccls{twocolumn}。

以下标题命令不受支持：\doccmd{subsubsection} 和 \doccmd{subparagraph}。




\chapter{故障排查与支持}
\label{ch:troubleshooting}

\section{\TL 网站}\label{sec:website}
\TL 项目的网站位于
\url{http://code.google.com/p/tufte-latex/}。在该网站上，你可以找到代码仓库、邮件列表、缺陷追踪器以及相关文档的链接。

\section{\TL 邮件列表}\label{sec:mailing-lists}
\TL 项目共有两个邮件列表：

\paragraph{讨论列表}
\texttt{tufte-latex} 讨论列表用于提问、寻求问题排查帮助，以及获取技术支持。版本发布公告也会发布在这个列表中。你可以通过以下地址订阅：\url{http://groups.google.com/group/tufte-latex}。

\paragraph{提交列表}
\texttt{tufte-latex-commits} 是一个只读列表。每当 \TL 代码被更新时，系统都会向该列表发送一封通知邮件。如果你希望及时了解代码的最新进展，可以订阅这个列表。订阅地址为：\url{http://groups.google.com/group/tufte-latex-commits}。

\section{获取帮助}\label{sec:getting-help}
如果你在使用某个 \TL 文档类时遇到问题，或者希望报告 bug、提出疑问，请向邮件列表发送邮件，或者访问项目网站。

为了帮助我们更快定位问题，请尽量使用 \docclsopt{debug} 类选项重新编译文档，并把生成的 \texttt{.log} 文件以及问题的简短描述一起发送到邮件列表。



\section{错误、警告与信息消息}\label{sec:tl-messages}
下面列出了 \TL 类生成的所有错误、警告和信息消息，并对其含义做了简要说明。
\index{error messages}\index{warning messages}\index{debug messages}

% Errors
\docmsg{错误：\doccmd{subparagraph} 在此文档类中未定义。}{%
\doccmd{subparagraph} 命令并未在 \TL 文档类中定义。如果你希望使用 \doccmd{subparagraph}，则需要自行重新定义它。有关 \TL 文档类中可用标题样式的说明，请参见第~\pageref{sec:headings} 页 “Headings” 一节。}

\docmsg{错误：\doccmd{subsubsection} 在此文档类中未定义。}{%
\doccmd{subsubsection} 命令并未在 \TL 文档类中定义。如果你希望使用 \doccmd{subsubsection}，则需要自行重新定义它。有关 \TL 文档类中可用标题样式的说明，请参见第~\pageref{sec:headings} 页 “Headings” 一节。}

\docmsg{错误：\doccmd{morefloats} 最多只能调用两次。请参见\par\noindent\ \ \ \ \ \ \ Tufte-LaTeX 文档中的其他解决办法。}{%
\LaTeX{} 初始会分配 18 个槽位存储浮动体。第一次调用 \doccmd{morefloats} 会额外分配 34 个槽位；第二次调用则再额外分配 26 个槽位。

\doccmd{morefloats} 最多只能调用两次。第三次调用时就会触发此错误消息。更多信息请参见第~\pageref{err:too-many-floats} 页。}

% Warnings
\docmsg{警告：选项 `\docopt{class option}' 不受支持，将被忽略。}{%
当你尝试在 \TL 文档类中使用不受支持的 \docopt{class option} 时，就会出现该警告。在这种情况下，该选项会被忽略。}

% Info / Debug messages
\docmsg{信息：`\docclsopt{symmetric}' 选项隐含启用 `\docclsopt{twoside}'。}{%
你指定了 \docclsopt{symmetric} 文档类选项。这个选项会自动同时启用 \docclsopt{twoside}。关于其详细说明，请参见第~\pageref{clsopt:symmetric} 页。}


\section{宏包依赖}\label{sec:dependencies}
下表列出了 \TL 文档类依赖的宏包。带星号者为可选依赖。
\begin{multicols}{2}
\begin{itemize}
  \item xifthen
  \item ifpdf*
  \item ifxetex*
  \item hyperref
  \item geometry
  \item ragged2e
  \item chngpage \emph{or} changepage
  \item paralist
  \item textcase
  \item soul*
  \item letterspace*
  \item setspace
  \item natbib \emph{and} bibentry
  \item optparams
  \item placeins
  \item mathpazo*
  \item helvet*
  \item fontenc
  \item beramono*
  \item fancyhdr
  \item xcolor
  \item textcomp
  \item titlesec
  \item titletoc
\end{itemize}
\end{multicols}




%%
% The back matter contains appendices, bibliographies, indices, glossaries, etc.








\backmatter

% \bibliography{sample-handout}
% \bibliographystyle{plainnat}


\chapter*{致谢}
\addcontentsline{toc}{chapter}{致谢}

本模板的版式与配色机制建立在若干开源工作之上,作者特此致谢:

\begin{itemize}
  \item \textbf{Tufte-LaTeX} 项目——本模板所基于的 \texttt{tufte-book} 文档类的来源,
    是整套书籍版式的排印基础。\\
    \url{https://github.com/Tufte-LaTeX/tufte-latex}
  \item \textbf{OUC Haide Beamer Template}——本书模板的配色机制与视觉语言的直接来源,
    \texttt{\textbackslash OUCSetPalette} 接口与 23 套配色均沿袭并扩展自该 beamer 模板。\\
    \url{https://github.com/Hatsuyuki017/OUC-Haide-Beamer-Template}
  \item \textbf{Mattia Puddu}——其 Tufte 风格书籍版式中受 VDQI 启发的标题页与目录页设计,
    直接启发了本模板的前置内容(front matter)排版。\\
    \url{https://mattiapuddu25.github.io/index.html}
\end{itemize}



\printindex

% Palette Introduction --- Chinese.
% Auto-generated by gen_palette_intro.py.  Do not edit by hand.
\chapter{调色板介绍}\label{ch:palette-intro}

本章按编号依次介绍 24 组主题调色板。每一组都保留三段固定结构——\emph{色彩哲学}阐明配色背后的世界观,\emph{六色配色组}列出 \texttt{pal1} 到 \texttt{pal6} 的色值与释义,随后的小节(\emph{设计思路}、\emph{色彩结构}或\emph{美学法则}等)再对每一色或每一条结构性原则展开说明。配色按编号顺序排列、自然分页,不再强制每页限定数量,只保证六色配色表保持完整、不被分割。

% ---------------------------------------------------------------------------
\bookpalettecardhead{0\enspace OUC Default}
\palettesubsec{色彩哲学}
\noindent 这是模板出厂时的底色。蓝是青年学者的颜色——从清晨第一页讲义、午后图书馆窗台,到深夜实验台的冷光,蓝构成学术工作的整个色温。六层由浅入深的蓝刻意维持理性的克制,模拟一本书从翻开到合上的全过程;它不是一组要被注视的色,而是一组让文字自身被看见的色。这是写作者与读者之间最不打扰的契约,也是中国海洋大学海德学院 LaTeX 模板希望被使用者长期相处的方式。

\begin{fullwidth}\small\noindent
\begin{tabularx}{\linewidth}{@{}clllX@{}}
\toprule\rowcolor{booktableheadcolor}
\booktableheadcell{} & \booktableheadcell{中文名} & \booktableheadcell{原名} & \booktableheadcell{RGB} & \booktableheadcell{释义} \\
\midrule
\bookcardswatch{239,246,255} & 纸海蓝 & Paper Foam & 239, 246, 255 & 翻开扉页时纸面反射的最浅一抹蓝,几近白 \\
\bookcardswatch{219,234,254} & 冰晶蓝 & Desert Foam & 219, 234, 254 & 实验室冷光下纸张投出的浅影,清洁而无负担 \\
\bookcardswatch{96,165,250} & 天玻璃蓝 & Sea Smoke & 96, 165, 250 & 窗外的天,是学者抬头喘息时的呼吸色 \\
\bookcardswatch{37,99,235} & 讲堂蓝 & Crisp Sky & 37, 99, 235 & 黑板上一笔粉笔勾出的强调,是论证之色 \\
\bookcardswatch{30,58,138} & 深思蓝 & Twilight Blue & 30, 58, 138 & 思考最深处的暗蓝,接近梦的边界 \\
\bookcardswatch{30,41,59} & 铅墨石 & Midnight Hull & 30, 41, 59 & 钢笔留在笔记本上的最后一行字 \\
\bottomrule
\end{tabularx}\end{fullwidth}

\palettesubsec{阅读弧线的色彩映射}
\begin{palettebullets}
\item \textbf{Paper Foam} --- 纸海蓝是「开卷」的呼吸口,几近于白,允许任何一种文字以最少的视觉负担进入读者的眼睛
\item \textbf{Frost Crystal} --- 冰晶蓝承担背景的次轻调,在长时间阅读中替代纯白,降低疲劳而不引起注意
\item \textbf{Sky Glass} --- 天玻璃蓝是六色中唯一带有「自然」气息的中调,把书页从全然的人工色阵列里拉回一次窗外的视野
\item \textbf{Lecture Sapphire} --- 讲堂蓝承担论证的强调职责,是六色中唯一允许「响」的色——但响得克制,像粉笔在板上敲下重音
\item \textbf{Reverie Navy} --- 深思蓝把思考压向最深处,适合做长引文、关键定理、章节扉页的暗调主色
\item \textbf{Slate Ink} --- 铅墨石几乎是黑,却保有蓝的内核,是正文文字与所有结构线条最终落笔的颜色——书在它的颜色里收束
\end{palettebullets}
\par\bigskip
% ---------------------------------------------------------------------------
\bookpalettecardhead{1\enspace Brunneophobia}
\palettesubsec{色彩哲学}
\noindent 名字直译为「恐褐症」,却用整组褐色直面这种偏见。褐不是廉价、不是过时,而是被现代审美刻意冷落的暖。从沙土、橘黄、铁锈到深焦,六色一路下沉,如同走进一间堆满旧书的木质工作室,鼻腔里有皮革、烟草、湿木的气息。色彩的勇气不在于鲜艳,而在于敢做被回避的那一种;这组配色把所有被嫌弃的褐重新收编进一种成熟的体面。

\begin{fullwidth}\small\noindent
\begin{tabularx}{\linewidth}{@{}clllX@{}}
\toprule\rowcolor{booktableheadcolor}
\booktableheadcell{} & \booktableheadcell{中文名} & \booktableheadcell{原名} & \booktableheadcell{RGB} & \booktableheadcell{释义} \\
\midrule
\bookcardswatch{238,211,180} & 沙浪米 & Lighthouse & 238, 211, 180 & 沙漠日出的最浅暖白,是褐色系的「呼吸口」 \\
\bookcardswatch{213,148,79} & 烟橘 & Sand Lily & 213, 148, 79 & 烟雾掩映下的橘黄,温暖却不张扬 \\
\bookcardswatch{213,148,79} & 蜂蜡金 & Reef Coral & 213, 148, 79 & 同色重复——表明这组褐刻意收紧色相 \\
\bookcardswatch{180,69,15} & 铁锈赤 & Tidal Mauve & 180, 69, 15 & 暴露在湿气中数十年的铁,带着血与土 \\
\bookcardswatch{86,67,53} & 旧皮褐 & Inkwell & 86, 67, 53 & 老沙发与皮装订书脊的颜色,有岁月气息 \\
\bookcardswatch{42,23,14} & 焦土黑 & Deep Tide & 42, 23, 14 & 篝火熄灭后炭灰里残留的最深褐 \\
\bottomrule
\end{tabularx}\end{fullwidth}

\palettesubsec{褐色家族的内部结构}
\begin{palettebullets}
\item \textbf{Desert Foam} --- 沙浪米作为整组配色的最浅调,承担所有「呼吸」与「留白」的职责,是褐色系里唯一不发声的色
\item \textbf{Smoke Tangerine} --- 烟橘负责中段的温度,把整组配色从严肃拉回日常,是这组色与读者之间最亲切的接触面
\item \textbf{Beeswax Honey} --- 蜂蜡金与烟橘共用色相,是整组配色刻意为之的「收紧」——意在证明克制本身即是力量
\item \textbf{Rust Crimson} --- 铁锈赤是这组色彩的「锋」——它把柔和的褐推向有侵略性的红,是唯一允许刺痛的色
\item \textbf{Aged Leather} --- 旧皮褐承担文献质感,是图书馆、档案、皮装订所共享的深调,把读物的物质感留在了色彩里
\item \textbf{Scorched Earth} --- 焦土黑是结构的根,所有上层暖色都因为它的存在而得以稳稳站立——克服恐褐症的最终一步,是承认黑的褐属性
\end{palettebullets}
\par\bigskip
% ---------------------------------------------------------------------------
\bookpalettecardhead{2\enspace Van Dyke}
\palettesubsec{色彩哲学}
\noindent 致敬十七世纪佛兰德画家 Anthony van Dyck——他的肖像画用「凡戴克褐」打底,以薄釉层层堆叠出贵族肌肤的透光感。这组配色从他的画布生长而来:粉调的肌肤、淡紫的丝绸、暗红的天鹅绒、深褐的背景,皆是巴洛克肖像的标准物质。六色之间没有戏剧冲突,只有褪色后的体面;它们组成的不是一幅画,而是一幅画被时间洗过之后,墙上仍然挂着的那种气场。

\begin{fullwidth}\small\noindent
\begin{tabularx}{\linewidth}{@{}clllX@{}}
\toprule\rowcolor{booktableheadcolor}
\booktableheadcell{} & \booktableheadcell{中文名} & \booktableheadcell{原名} & \booktableheadcell{RGB} & \booktableheadcell{释义} \\
\midrule
\bookcardswatch{236,194,188} & 肌肤粉 & Vleeskleur & 236, 194, 188 & 巴洛克肖像中贵族面颊的薄釉粉调 \\
\bookcardswatch{169,159,191} & 丝绸藤 & Zijden Sering & 169, 159, 191 & 折光丝绸在阴影里的淡紫,优雅而冷 \\
\bookcardswatch{169,159,191} & 烟雾紫 & Rooklila & 169, 159, 191 & 同前——丝绸两面在画布上几乎不分 \\
\bookcardswatch{191,113,133} & 玫瑰呢 & Rozenfluweel & 191, 113, 133 & 天鹅绒礼服里渗出的玫瑰红,沉而不艳 \\
\bookcardswatch{68,60,94} & 暗紫绒 & Pruimenpurper & 68, 60, 94 & 背景帷幕的深紫,把光衬得更亮 \\
\bookcardswatch{61,43,39} & 凡戴克褐 & Van Dijck-bruin & 61, 43, 39 & 该画家的标志底色,以沥青颜料调和而成 \\
\bottomrule
\end{tabularx}\end{fullwidth}

\palettesubsec{肖像深度的色彩映射}
\begin{palettebullets}
\item \textbf{Carnation Skin} --- 肌肤粉是画面最前景——画中人物的面颊、手指、颈侧,是观看者第一眼接触到的「人」
\item \textbf{Silken Wisteria} --- 丝绸藤是衣料的高光,把柔软的折光感留在了色相中,优雅而带着距离
\item \textbf{Smoke Lilac} --- 烟雾紫与丝绸藤同色,刻意制造重复,呼应巴洛克绘画里对单色叠层的偏爱
\item \textbf{Rose Velvet} --- 玫瑰呢是中景礼服的颜色,把人物从画中抬起、推近,是这组配色的色温支点
\item \textbf{Plum Plush} --- 暗紫绒是背景帷幕,负责把前景的人物从黑暗中托出,自己几乎不被注意
\item \textbf{Van Dyke Brown} --- 凡戴克褐是整幅画的「地」——是颜料、画布、画框、墙面共同沉淀下来的底色,任何高光都因它而立
\end{palettebullets}
\par\bigskip
% ---------------------------------------------------------------------------
\bookpalettecardhead{3\enspace Back in Black}
\palettesubsec{色彩哲学}
\noindent 致敬 AC/DC 1980 年的同名专辑,但配色本身比摇滚更克制。粉与灰柔化了黑的攻击性,使整组色更接近后朋克美学:皮夹克、烟雾、舞台灯熄灭后那一刻的余温。粉是疤痕,黑是夜——它们之间所有的灰,是青春过后还留在身体里的回声。这是一组「成熟版的叛逆」,把锋芒收进了体面之内。

\begin{fullwidth}\small\noindent
\begin{tabularx}{\linewidth}{@{}clllX@{}}
\toprule\rowcolor{booktableheadcolor}
\booktableheadcell{} & \booktableheadcell{中文名} & \booktableheadcell{原名} & \booktableheadcell{RGB} & \booktableheadcell{释义} \\
\midrule
\bookcardswatch{240,217,228} & 婴儿粉 & Powder Pink & 240, 217, 228 & 黑暗里唯一不设防的颜色,是叛逆的反面 \\
\bookcardswatch{193,160,172} & 灰玫瑰 & Dust Rose & 193, 160, 172 & 蒙了一层烟的粉,失去鲜艳后的体面 \\
\bookcardswatch{193,160,172} & 烬粉 & Ember Pink & 193, 160, 172 & 重复以加强中调,粉的两种语义并置 \\
\bookcardswatch{128,108,121} & 雾紫灰 & Smoky Mauve & 128, 108, 121 & 舞台灯下烟雾的灰紫,既不明也不暗 \\
\bookcardswatch{74,63,75} & 黑天鹅绒 & Black Velvet & 74, 63, 75 & 皮夹克的暗影,介于黑与紫之间 \\
\bookcardswatch{22,19,21} & 夜墨 & Inkwell Black & 22, 19, 21 & 一组配色最沉的锚点,几乎黑 \\
\bottomrule
\end{tabularx}\end{fullwidth}

\palettesubsec{后朋克美学的色彩映射}
\begin{palettebullets}
\item \textbf{Powder Pink} --- 婴儿粉是这组配色的「不设防」——它是黑底之上唯一一个温柔的入口,叛逆的反面才是叛逆的内核
\item \textbf{Dust Rose} --- 灰玫瑰把粉的鲜艳压了下去,是青春之后留下的余韵,优雅地承认了「不再年轻」
\item \textbf{Ember Pink} --- 烬粉与灰玫瑰共享色相,是「重复即强调」的视觉修辞,把中调的灰玫瑰提升为整组色的真正核心
\item \textbf{Smoky Mauve} --- 雾紫灰是舞台烟雾的颜色,过渡了粉与黑,使两端的对立看起来像同一种情绪的两面
\item \textbf{Black Velvet} --- 黑天鹅绒是皮夹克的暗影,带着紫的呼吸——黑不必是死的,它可以是一种被穿在身上的物质
\item \textbf{Inkwell Black} --- 夜墨是最沉的锚点,接近纯黑却保有微紫底调,是整组色的句号——但句号之后还可以有下一首歌
\end{palettebullets}
\par\bigskip
% ---------------------------------------------------------------------------
\bookpalettecardhead{4\enspace Belle of the Ball}
\palettesubsec{色彩哲学}
\noindent 舞会上那位最被注视的女子从不穿原色——她穿的是被无数烛光暖化过的旧色:陶土、铜锈、橄榄、绛红。这组配色拒绝二十世纪的合成染料,从十九世纪欧洲晚宴厅与殖民地热带气候的混合美学中提取——壁炉、桃花心木、暮色窗帘、染了酒渍的桌布。被注视并不靠鲜艳,而靠所有颜色都已被时间柔化过的那种「被生活过的高级感」。

\begin{fullwidth}\small\noindent
\begin{tabularx}{\linewidth}{@{}clllX@{}}
\toprule\rowcolor{booktableheadcolor}
\booktableheadcell{} & \booktableheadcell{中文名} & \booktableheadcell{原名} & \booktableheadcell{RGB} & \booktableheadcell{释义} \\
\midrule
\bookcardswatch{226,203,192} & 香槟珠光 & Champagne Pearl & 226, 203, 192 & 舞会大厅烛光下女子肌肤的珠光底色 \\
\bookcardswatch{206,171,150} & 玫瑰陶土 & Rose Terracotta & 206, 171, 150 & 复古口红与陶土色的中间调,被注视的暖色 \\
\bookcardswatch{210,135,106} & 陶釉橘 & Glazed Clay & 210, 135, 106 & 陶器釉面的橘红,温润而成熟 \\
\bookcardswatch{229,74,57} & 辣椒朱 & Chilli Scarlet & 229, 74, 57 & 热带气候里晾在阳光下的辣椒色,是这组配色的「亮点」 \\
\bookcardswatch{118,118,44} & 橄榄绿 & Olive Verdigris & 118, 118, 44 & 橄榄叶被氧化后的中性绿,带些铜锈 \\
\bookcardswatch{53,77,4} & 苔藓深绿 & Moss Deepwood & 53, 77, 4 & 桃花心木家具与深色窗帘的森林暗调 \\
\bottomrule
\end{tabularx}\end{fullwidth}

\palettesubsec{被注视者的色彩策略}
\begin{palettebullets}
\item \textbf{Champagne Pearl} --- 香槟珠光是被烛光暖化的肌肤底色,是整组配色的「呼吸」,也是被注视者最不设防的一层
\item \textbf{Rose Terracotta} --- 玫瑰陶土承担「复古性感」,是被时间筛过的口红色,既是装饰也是身份的证明
\item \textbf{Glazed Clay} --- 陶釉橘是手作之色,把工业感从配色中排除,留下手心温度与陶土的物质感
\item \textbf{Chilli Scarlet} --- 辣椒朱是全组唯一允许「亮」的色,是裙摆的红、戒指的红、唇上最后一笔的红——但只能小面积出现
\item \textbf{Olive Verdigris} --- 橄榄绿是这组色彩的冷调支点,把所有热带暖色拉回欧洲晚宴厅的克制,是配色的色温调节器
\item \textbf{Moss Deepwood} --- 苔藓深绿是背景与家具的暗调,是这组配色的「重」,没有它,所有的亮都会失去支点
\end{palettebullets}
\par\bigskip
% ---------------------------------------------------------------------------
\bookpalettecardhead{5\enspace Pine Tree}
\palettesubsec{色彩哲学}
\noindent 以一棵北方松树作为坐标,从树冠到树根逐层收色:阳光打在新生针叶上的暖金、被风吹松动的旧针、午后树荫里的玫瑰陶土、树干裂痕里渗出的赭色,以及最底层埋在腐殖土里的森林暗墨。配色逻辑刻意离开「绿色」——一棵真正的松树在视觉上远比绿色复杂得多,它的颜色其实是阳光、岁月、土壤、伤口与影子共同写下的一段年表。

\begin{fullwidth}\small\noindent
\begin{tabularx}{\linewidth}{@{}clllX@{}}
\toprule\rowcolor{booktableheadcolor}
\booktableheadcell{} & \booktableheadcell{中文名} & \booktableheadcell{原名} & \booktableheadcell{RGB} & \booktableheadcell{释义} \\
\midrule
\bookcardswatch{238,200,111} & 阳针金 & Sunlit Needle & 238, 200, 111 & 阳光打在新生针叶上的暖金 \\
\bookcardswatch{222,166,32} & 蜂蜜松脂 & Honey Resin & 222, 166, 32 & 树皮裂口渗出的松脂金,是松树自愈的颜色 \\
\bookcardswatch{222,166,32} & 陈年金 & Aged Gold & 222, 166, 32 & 同色再现,使中调金更稳 \\
\bookcardswatch{177,120,133} & 幽影玫瑰 & Shadow Rose & 177, 120, 133 & 午后树荫里的玫瑰陶土,带紫调 \\
\bookcardswatch{167,88,26} & 树皮赭 & Bark Ochre & 167, 88, 26 & 树干主体的赭褐,是松树最诚实的颜色 \\
\bookcardswatch{43,47,34} & 林墨 & Forest Ink & 43, 47, 34 & 腐殖土与树根处的森林暗墨 \\
\bottomrule
\end{tabularx}\end{fullwidth}

\palettesubsec{松树年表的色彩映射}
\begin{palettebullets}
\item \textbf{Sunlit Needle} --- 阳针金是松树的「今天」——新生的、被光选中的、还没有故事的那一段时间
\item \textbf{Honey Resin} --- 蜂蜜松脂是松树的「伤口」——树脂封住裂痕,使损伤变成琥珀,是这组色彩最具时间感的一色
\item \textbf{Aged Gold} --- 陈年金与蜂蜜松脂共享色相,是中调金的「再说一遍」,把时间在颜色里复读一次
\item \textbf{Shadow Rose} --- 幽影玫瑰是松树投在地面的紫红色阴影,是树与土壤之间的呼吸界面
\item \textbf{Bark Ochre} --- 树皮赭是树干本体,是最不需要解释的色——它的赭褐就是松树这个生命体的物质本色
\item \textbf{Forest Ink} --- 林墨是松树的「根之色」,沉在腐殖土里,接近黑而仍属绿——它告诉我们,一棵活树的根,从不真正与黑相同
\end{palettebullets}
\par\bigskip
% ---------------------------------------------------------------------------
\bookpalettecardhead{6\enspace Provence Blue}
\palettesubsec{色彩哲学}
\noindent 普罗旺斯的蓝从不属于海——它属于薰衣草田尽头那片被夕阳压低饱和的蓝灰,属于午后石灰岩房屋窗框上经年褪色的孔雀蓝,属于一户人家晾在院里的旧亚麻床单与铁锈门把手共同呼吸的色温。这组配色拒绝地中海明信片式的浪漫,选择更内陆、更乡村、更被时间柔化过的法国南部蓝调——它的浪漫不在颜色的鲜艳里,而在颜色被使用过后留下的痕迹里。

\begin{fullwidth}\small\noindent
\begin{tabularx}{\linewidth}{@{}clllX@{}}
\toprule\rowcolor{booktableheadcolor}
\booktableheadcell{} & \booktableheadcell{中文名} & \booktableheadcell{原名} & \booktableheadcell{RGB} & \booktableheadcell{释义} \\
\midrule
\bookcardswatch{170,188,175} & 橄榄银灰 & Argent d'Olivier & 170, 188, 175 & 橄榄树叶背面的银灰,带绿调 \\
\bookcardswatch{137,156,154} & 雾镜蓝 & Miroir de Brume & 137, 156, 154 & 清晨雾气压在水面上的灰蓝 \\
\bookcardswatch{137,156,154} & 宿雨青 & Pluie Pass\'ee & 137, 156, 154 & 同色再现,巩固中调 \\
\bookcardswatch{110,124,139} & 石窗蓝 & Bleu de Volet & 110, 124, 139 & 石灰岩房屋窗框上的孔雀蓝褪色后的色 \\
\bookcardswatch{82,92,121} & 薰衣草夜 & Cr\'epuscule Lavande & 82, 92, 121 & 薰衣草田尽头被夕阳压低的暗紫蓝 \\
\bookcardswatch{53,66,94} & 橄榄夜空 & Nuit d'Olivier & 53, 66, 94 & 橄榄园上方夜空的深蓝,带土壤气息 \\
\bottomrule
\end{tabularx}\end{fullwidth}

\palettesubsec{南法乡村蓝调的色彩结构}
\begin{palettebullets}
\item \textbf{Olive Silver} --- 橄榄银灰是这组色彩里唯一的绿,是橄榄树叶背面在风里翻起的银光,是普罗旺斯空气的物质化
\item \textbf{Mist Mirror} --- 雾镜蓝是清晨水面的镜像,把蓝从天空借回地面,呼吸感最强的中调
\item \textbf{Spent Rain} --- 宿雨青与雾镜蓝同色,是「雨后第二日」的色温,刻意重复以加强配色的湿度
\item \textbf{Window Stone Blue} --- 石窗蓝把这组色彩从自然带回人造——它是被无数次刷过、又被无数次淡去的人手痕迹
\item \textbf{Lavender Dusk} --- 薰衣草夜是配色的灵魂——它把薰衣草从「紫」推向「暗紫蓝」,这一步,正是普罗旺斯之所以为普罗旺斯
\item \textbf{Olive Night} --- 橄榄夜空是这组色彩的根,把白日的所有蓝色收进了一片夜空,接近黑而仍是蓝
\end{palettebullets}
\par\bigskip
% ---------------------------------------------------------------------------
\bookpalettecardhead{7\enspace Fresco Blue}
\palettesubsec{色彩哲学}
\noindent 来自意大利文艺复兴时期教堂壁画的天空——那种被湿壁画(Fresco)技法吸入灰泥里、再被五百年烛烟氧化过的青蓝。它不是天空真实的颜色,而是匠人用青金石、孔雀石与石灰浆调出的「神圣的天空」,介于人间与天国之间。这组配色从最浅的天到最深的夜按梯度排列,模拟壁画从穹顶向地面跌落时蓝色的层层加深——它是一种被建筑收容过的天。

\begin{fullwidth}\small\noindent
\begin{tabularx}{\linewidth}{@{}clllX@{}}
\toprule\rowcolor{booktableheadcolor}
\booktableheadcell{} & \booktableheadcell{中文名} & \booktableheadcell{原名} & \booktableheadcell{RGB} & \booktableheadcell{释义} \\
\midrule
\bookcardswatch{166,224,244} & 穹顶浅天 & Cielo di Volta & 166, 224, 244 & 教堂穹顶最高处壁画的浅蓝 \\
\bookcardswatch{71,169,207} & 钴釉蓝 & Vetrina Cobalto & 71, 169, 207 & 青金石与石灰浆调和后的中调蓝 \\
\bookcardswatch{71,169,207} & 青金蓝 & Lapislazzuli & 71, 169, 207 & 同色再现,巩固壁画的中段 \\
\bookcardswatch{9,121,158} & 深殿蓝 & Azzurro d'Abside & 9, 121, 158 & 教堂半圆殿背景里五百年未褪的深蓝 \\
\bookcardswatch{4,75,102} & 夜祷蓝 & Azzurro della Veglia & 4, 75, 102 & 夜祷时灯火映在壁画下半部分的暗蓝 \\
\bookcardswatch{2,31,46} & 地宫蓝 & Azzurro della Cripta & 2, 31, 46 & 地宫与暗廊里壁画几近隐没的最深蓝 \\
\bottomrule
\end{tabularx}\end{fullwidth}

\palettesubsec{壁画天空的垂直色谱}
\begin{palettebullets}
\item \textbf{Vault Sky} --- 穹顶浅天是壁画最高处,光线打在灰泥上反射回来的浅蓝,是这组配色的「天」
\item \textbf{Cobalt Glaze} --- 钴釉蓝是中段的主调,是教堂内信众抬头时看见的那一抹饱满
\item \textbf{Lapis Renewed} --- 青金蓝与钴釉蓝同色,是「再次确认」——壁画的蓝从不依靠一次涂抹,它依靠层层叠加
\item \textbf{Apse Blue} --- 深殿蓝是教堂建筑里光线最难抵达的深处,壁画在那里凝固成不褪色的深蓝
\item \textbf{Vigil Blue} --- 夜祷蓝是夜晚的色温——同一面壁画,在油灯下显现出一种与白日截然不同的暗蓝
\item \textbf{Crypt Blue} --- 地宫蓝是壁画几乎被黑暗吞没的最深处,几近黑而仍属蓝——这是一组色彩中最沉默的「相信」
\end{palettebullets}
\par\bigskip
% ---------------------------------------------------------------------------
\bookpalettecardhead{8\enspace Monet}
\palettesubsec{色彩哲学}
\noindent 印象派的核心命题是:光本身有颜色。莫奈一生反复画同一座干草堆、同一座大教堂、同一池睡莲,只为捕捉「同一个对象在不同时刻拥有不同色彩」的事实。这组配色不复刻任何一张莫奈作品,而是从他全部画作里抽出最具印象派身份的六个色相:奶油画布、玫瑰干草、绿草地、夏池水、晨雾蓝、深苇影——它们之间的关系不是构图,而是同一束光在不同物质上的折射。

\begin{fullwidth}\small\noindent
\begin{tabularx}{\linewidth}{@{}clllX@{}}
\toprule\rowcolor{booktableheadcolor}
\booktableheadcell{} & \booktableheadcell{中文名} & \booktableheadcell{原名} & \booktableheadcell{RGB} & \booktableheadcell{释义} \\
\midrule
\bookcardswatch{247,244,213} & 画布奶油 & Toile Cr\`eme & 247, 244, 213 & 莫奈打底未上色画布的奶油底 \\
\bookcardswatch{211,150,140} & 玫瑰干草 & Foin Rose & 211, 150, 140 & 《干草堆》系列中黄昏时干草顶端的玫瑰 \\
\bookcardswatch{211,150,140} & 夕辉粉 & Lueur du Soir & 211, 150, 140 & 同色再现,延长黄昏的色温 \\
\bookcardswatch{131,153,88} & 草地绿 & Pr\'e Vert & 131, 153, 88 & 吉维尼花园午后草地的中调绿 \\
\bookcardswatch{16,86,102} & 池水蓝 & \'Etang Aigue & 16, 86, 102 & 《睡莲》系列水面下倒映的深蓝 \\
\bookcardswatch{10,51,35} & 芦苇暗影 & Ombre des Roseaux & 10, 51, 35 & 池边芦苇丛根部的深绿,接近墨 \\
\bottomrule
\end{tabularx}\end{fullwidth}

\palettesubsec{印象派光色的色彩结构}
\begin{palettebullets}
\item \textbf{Canvas Cream} --- 画布奶油是这组配色的「未完成」——是莫奈每一次起笔之前,画布本身的呼吸
\item \textbf{Hay Rose} --- 玫瑰干草是印象派最具识别性的一笔——它把「黄色」推向「玫瑰」,正是因为黄昏的光改写了一切固有色
\item \textbf{Evening Glow} --- 夕辉粉与玫瑰干草同色,是「延长黄昏」——印象派之所以反复画同一对象,正为多次捕捉这条色
\item \textbf{Meadow Green} --- 草地绿是配色的中调,是阳光打在新草上时的真实绿,既不冷也不暖
\item \textbf{Pond Aqua} --- 池水蓝是莫奈晚年《睡莲》中最具识别性的色相,把蓝从天空借给水,是这组配色的情绪转折点
\item \textbf{Reed Shadow} --- 芦苇暗影是这组配色的「根」,几近墨而仍是绿——印象派的暗部从不真正接受黑
\end{palettebullets}
\par\bigskip
% ---------------------------------------------------------------------------
\bookpalettecardhead{9\enspace Narcissus}
\palettesubsec{色彩哲学}
\noindent 不是希腊神话里的水仙,而是植物学意义上的水仙——一种从球根抽出花茎、在冬末春初开放的花。它的颜色谱系从苍白的鳞茎到金色的副冠、再到深褐的泥土根系,是一段「从地下到地面」的物质叙事。这组配色刻意去掉「绿叶」这一最显而易见的项,从而把视线集中在花本身从孕育、开放到回归大地的全过程,是一组以「时间纵轴」为结构的植物色谱。

\begin{fullwidth}\small\noindent
\begin{tabularx}{\linewidth}{@{}clllX@{}}
\toprule\rowcolor{booktableheadcolor}
\booktableheadcell{} & \booktableheadcell{中文名} & \booktableheadcell{原名} & \booktableheadcell{RGB} & \booktableheadcell{释义} \\
\midrule
\bookcardswatch{221,213,200} & 鳞茎乳 & Bulbus Lacteus & 221, 213, 200 & 刚切开的水仙球根截面的乳白 \\
\bookcardswatch{185,149,144} & 花瓣粉褐 & Petalum Roseobrunneum & 185, 149, 144 & 落瓣压在泥土里的灰玫瑰 \\
\bookcardswatch{185,149,144} & 退色玫瑰 & Rosa Defloruit & 185, 149, 144 & 同色再现,代表「再用一次」的语义 \\
\bookcardswatch{199,149,72} & 水仙金 & Aurum Narcissi & 199, 149, 72 & 副冠中央的暖金,是配色的标识色 \\
\bookcardswatch{190,108,26} & 蜂蜡铜 & Aes Mellitum & 190, 108, 26 & 球根外壳氧化后的橙褐,介于土与花之间 \\
\bookcardswatch{110,60,31} & 泥土深 & Humus Profundus & 110, 60, 31 & 水仙根系所在的腐殖土深褐,是根之色 \\
\bottomrule
\end{tabularx}\end{fullwidth}

\palettesubsec{水仙生命周期的色彩切片}
\begin{palettebullets}
\item \textbf{Bulb Milk} --- 鳞茎乳是水仙的「未开」,是植物在尚未成为花之前已经拥有的全部潜在颜色
\item \textbf{Petal Brown-Rose} --- 花瓣粉褐是花的「凋落」,花瓣离开花茎后被泥土与雨水共同改写的颜色
\item \textbf{Faded Rose} --- 退色玫瑰与花瓣粉褐同色,是「记忆中的同一片花瓣」——同一种凋落,在记忆里出现了第二次
\item \textbf{Narcissus Gold} --- 水仙金是这组配色的核心,是副冠中央那一抹只在春天初醒时出现的暖金,是「开放」本身的颜色
\item \textbf{Honey Bronze} --- 蜂蜡铜把花的金推向了土的赭,是花朵从开放走向凋落、从空中回到地下的过渡色
\item \textbf{Soil Deep} --- 泥土深是水仙的「来处」与「归处」,几近黑而保暖——这一组色彩告诉我们:任何盛开都从黑暗里来,也回去
\end{palettebullets}
\par\bigskip
% ---------------------------------------------------------------------------
\bookpalettecardhead{10\enspace Roman Empire}
\palettesubsec{色彩哲学}
\noindent 以罗马帝国的物质遗产为坐标:大理石、军团战袍、元老院托加、月桂冠、凯旋金、骨螺紫。六色不是设计而来,而是从罗马城本身长出来的——神庙的石、战场的血、广场的权、帝王的纹,共同构成一组配色的政治考古学。每一色都有可追溯的物质起源。

\begin{fullwidth}\small\noindent
\begin{tabularx}{\linewidth}{@{}clllX@{}}
\toprule\rowcolor{booktableheadcolor}
\booktableheadcell{} & \booktableheadcell{中文名} & \booktableheadcell{原名} & \booktableheadcell{RGB} & \booktableheadcell{释义} \\
\midrule
\bookcardswatch{236,232,225} & 大理石 & Marmor Carrariense & 236, 232, 225 & 卡拉拉大理石的冷白,略带暖灰,是神庙与雕像的物质 \\
\bookcardswatch{212,175,55} & 荣耀金 & Aurum Gloriae & 212, 175, 55 & 凯旋式黄金的中性金,介于铸币与神像镀金之间 \\
\bookcardswatch{74,110,65} & 桂冠 & Laurus Viridis & 74, 110, 65 & 月桂冠叶片的哑光深绿,取干燥叶片的沉稳色调 \\
\bookcardswatch{180,30,30} & 罗马军团 & Sagum Legionarium & 180, 30, 30 & 军团战袍(Sagum)的鲜烈战红,饱和而有力,象征铁与血 \\
\bookcardswatch{120,20,40} & 元老院 & Clavus Senatorius & 120, 20, 40 & 元老院托加袍缘的深沉酒红(Clavus),比军团红更内敛、更权贵 \\
\bookcardswatch{88,28,90} & 奥古斯都紫 & Purpura Tyria & 88, 28, 90 & 骨螺紫染料的历史色,价比黄金,专属皇权 \\
\bottomrule
\end{tabularx}\end{fullwidth}

\palettesubsec{设计思路}
\begin{palettebullets}
\item \textbf{Carrara Marble} --- 取自卡拉拉大理石的冷白底色,略带暖灰,呈现神庙与雕像的质感
\item \textbf{Legion Crimson} --- 军团战袍(Sagum)的鲜烈战红,饱和而有力,象征铁与血
\item \textbf{Senate Bordeaux} --- 元老院托加袍缘的深沉酒红(Clavus),比军团红更内敛、更权贵
\item \textbf{Laurel Viridis} --- 月桂冠叶片的哑光深绿,非翠绿,取干燥叶片的沉稳色调
\item \textbf{Gloria Aurum} --- 凯旋式黄金的中性金,介于铸币与神像镀金之间,不过暖也不过冷
\item \textbf{Tyrian Purple} --- 骨螺紫染料的历史色,价比黄金,专属皇权,深邃而神秘
\end{palettebullets}
\par\bigskip
% ---------------------------------------------------------------------------
\bookpalettecardhead{11\enspace Greece}
\palettesubsec{色彩哲学}
\noindent 以大理石白为底,民主蓝统领视觉重心,金色做高光点缀;陶器赤与橄榄绿构成一对冷暖平衡的中间调,葡萄紫作为最深的暗部锚点。大理石选用帕罗斯岛(Paros)产的暖白,比罗马卡拉拉更偏象牙,是帕特农神庙的实际用料;民主蓝取爱琴海深水区的中性蔚蓝,呼应雅典公民在广场(Agora)凝望的天与海。整体在地中海阳光下呈现出既庄重又充满生命力的古典张力。

\begin{fullwidth}\small\noindent
\begin{tabularx}{\linewidth}{@{}clllX@{}}
\toprule\rowcolor{booktableheadcolor}
\booktableheadcell{} & \booktableheadcell{中文名} & \booktableheadcell{原名} & \booktableheadcell{RGB} & \booktableheadcell{释义} \\
\midrule
\bookcardswatch{245,240,228} & 帕罗斯大理石 & \foreignname{Παριανός Μάρμαρος} \textperiodcentered\ Parianos Marmaros & 245, 240, 228 & 帕特农神庙的实际用料,比卡拉拉更偏象牙的暖白 \\
\bookcardswatch{212,175,55} & 荣耀金 & \foreignname{Χρυσός} \textperiodcentered\ Chrysos & 212, 175, 55 & 与罗马同源的中性金,呼应神明与宗教仪式的高光 \\
\bookcardswatch{48,105,175} & 民主蓝 & \foreignname{Κυανός Αγοράς} \textperiodcentered\ Kyanos Agoras & 48, 105, 175 & 爱琴海深水区的中性蔚蓝,雅典公民在广场凝望的天与海 \\
\bookcardswatch{188,82,38} & 陶器赤 & \foreignname{Αττικός Πηλός} \textperiodcentered\ Attikos Pelos & 188, 82, 38 & 红绘式陶器的底色,希腊人记录神话与英雄史诗最具标志性的媒介色 \\
\bookcardswatch{98,128,48} & 橄榄银绿 & \foreignname{Ελιά Αθηνάς} \textperiodcentered\ Elia Athenas & 98, 128, 48 & 雅典娜赠予城邦的橄榄树,和平与智慧的颜色 \\
\bookcardswatch{90,42,92} & 狄俄尼索斯葡萄紫 & \foreignname{Διονυσιακή Σταφυλή} \textperiodcentered\ Dionysiake Staphyle & 90, 42, 92 & 酿酒葡萄的深邃果紫,酒神、戏剧与狂欢节的灵魂色 \\
\bottomrule
\end{tabularx}\end{fullwidth}

\palettesubsec{配色逻辑}
\begin{palettebullets}
\item \textbf{Parian Marble} --- 大理石白作为底色,是希腊建筑与雕塑的物质本相——比罗马卡拉拉更暖、更象牙,带着地中海阳光过滤后的体温
\item \textbf{Agora Kyanos} --- 民主蓝统领整组色彩的视觉重心,既是爱琴海的天与海,也是雅典公民在广场上看见的「公共空间」的颜色
\item \textbf{Gloria Aurum} --- 荣耀金作为点缀,以最小面积承担神圣与神明的视觉重量,与罗马同源却用法不同——希腊的金更克制
\item \textbf{Attic Terracotta} --- 陶器赤是希腊的「叙事色」——红绘陶器把神话与英雄史诗压在这一色里,是希腊文化最具识别性的媒介
\item \textbf{Athena's Olive} --- 橄榄银绿与陶器赤构成冷暖平衡的中间调,是「智慧」与「叙事」的色彩对位
\item \textbf{Dionysian Grape} --- 葡萄紫是整组配色最深的暗部锚点,是酒神与戏剧的颜色——希腊文化的另一极:神圣之下的狂欢
\end{palettebullets}
\par\bigskip
% ---------------------------------------------------------------------------
\bookpalettecardhead{12\enspace Kanagawa}
\palettesubsec{色彩哲学}
\noindent 《神奈川冲浪里》的革命性在于葛饰北斎将当时刚传入日本的普鲁士蓝(ベロ藍)引入浮世绘,以单一色系的深浅变奏构建出整幅画面的戏剧张力,配合墨线轮廓与富士山的远景留白,形成极具压缩感的视觉美学。这组配色以普鲁士蓝家族为核心,从墨黑$\to$深海蓝$\to$浪蓝$\to$富士霞$\to$暮天$\to$波白构成一条完整的明度递进链,如同浪从深处涌起、在浪尖碎裂成雪沫的瞬间。

\begin{fullwidth}\small\noindent
\begin{tabularx}{\linewidth}{@{}clllX@{}}
\toprule\rowcolor{booktableheadcolor}
\booktableheadcell{} & \booktableheadcell{中文名} & \booktableheadcell{原名} & \booktableheadcell{RGB} & \booktableheadcell{释义} \\
\midrule
\bookcardswatch{237,233,222} & 波白 & \foreignname{波白} \textperiodcentered\ Nami-shiro & 237, 233, 222 & 浪尖破碎的泡沫,非纯白,带和纸的微暖底色 \\
\bookcardswatch{208,224,238} & 暮天 & \foreignname{暮天} \textperiodcentered\ Boten & 208, 224, 238 & 海平线处的天空渐变色,普鲁士蓝淡化至苍茫 \\
\bookcardswatch{150,186,210} & 富士霞 & \foreignname{富士霞} \textperiodcentered\ Fuji-gasumi & 150, 186, 210 & 远景富士山的淡蓝轮廓,虚化而沉静 \\
\bookcardswatch{26,78,132} & 普鲁士蓝 & \foreignname{ベロ藍} \textperiodcentered\ Bero-ai & 26, 78, 132 & 画面主体——普鲁士蓝浪身,江户最具标志性的舶来颜料 \\
\bookcardswatch{13,38,76} & 深海 & \foreignname{深海} \textperiodcentered\ Shinkai & 13, 38, 76 & 浪底最深处的墨蓝,近乎黑而仍保有蓝的冷意 \\
\bookcardswatch{29,25,35} & 墨 & \foreignname{墨} \textperiodcentered\ Sumi & 29, 25, 35 & 木版印刷的墨线,统摄全图骨架,非纯黑而含深紫底调 \\
\bottomrule
\end{tabularx}\end{fullwidth}

\palettesubsec{配色逻辑}
\begin{palettebullets}
\item \textbf{Nami-shiro} --- 波白是浪尖那一瞬间的「破碎」,带和纸的微暖底色,是整组色彩唯一的高光
\item \textbf{Boten} --- 暮天是海平线处天空的过渡色,把普鲁士蓝淡化至几乎消失,是画面「呼吸」的色
\item \textbf{Fuji-gasumi} --- 富士霞是远景富士山的淡蓝,是整幅画唯一的「退隐」,把焦点推回浪本身
\item \textbf{Bero-ai} --- 普鲁士蓝是浪身的主色,是这幅画的革命所在——一种单一外来颜料,撑起了整幅作品的戏剧张力
\item \textbf{Shinkai} --- 深海是浪底最深处的墨蓝,是「压迫感」与「深渊」的色彩载体,近乎黑而仍是蓝
\item \textbf{Sumi} --- 墨是木版印刷的骨架,统摄全图,是这组配色的「形」,没有它,六色蓝无可附着
\end{palettebullets}
\par\bigskip
% ---------------------------------------------------------------------------
\bookpalettecardhead{13\enspace Starry Night}
\palettesubsec{色彩哲学}
\noindent 梵高的《星月夜》并非印象派而是后印象派的巅峰之作——他不再捕捉光的瞬间,而是将内心的情绪直接燃烧进笔触。画面以钴蓝与群青为骨,用漩涡状的笔触赋予夜空以生命,月光与星芒的铬黄在蓝色的重压下迸裂成光爆,柏树如黑色火焰直刺苍穹,村庄的微光则是整幅画中唯一属于人间的暖。三层蓝构成一条由深至浅的情绪纵深轴,模拟梵高螺旋笔触的叠压节奏。

\begin{fullwidth}\small\noindent
\begin{tabularx}{\linewidth}{@{}clllX@{}}
\toprule\rowcolor{booktableheadcolor}
\booktableheadcell{} & \booktableheadcell{中文名} & \booktableheadcell{原名} & \booktableheadcell{RGB} & \booktableheadcell{释义} \\
\midrule
\bookcardswatch{240,208,68} & 月华铬黄 & Lumi\`ere Lunaire & 240, 208, 68 & 月亮与星芒的爆裂光晕,梵高常用铬黄表达极致的生命张力 \\
\bookcardswatch{198,140,52} & 村灯琥珀 & Lueurs du Village & 198, 140, 52 & 村庄窗口透出的暖黄,与冰冷蓝色星空形成全画唯一的冷暖对峙 \\
\bookcardswatch{105,155,200} & 晨曦苍蓝 & Aube Glac\'ee & 105, 155, 200 & 漩涡边缘的淡蓝过渡,深夜与月光交界处最微妙的呼吸感 \\
\bookcardswatch{48,96,165} & 星涡群青 & Tourbillon Outremer & 48, 96, 165 & 漩涡主体的群青蓝,后印象派用色的标志——非自然色而是情绪色 \\
\bookcardswatch{22,50,30} & 柏影墨绿 & Cypr\`es Nocturne & 22, 50, 30 & 前景柏树的深邃暗绿,几近于黑,是画面中唯一「拒绝光」的存在 \\
\bookcardswatch{20,36,88} & 深夜钴蓝 & Minuit Cobalt & 20, 36, 88 & 夜空最深处的底色,饱含压迫性的情绪重量,梵高精神世界的深渊底色 \\
\bottomrule
\end{tabularx}\end{fullwidth}

\palettesubsec{配色逻辑}
\begin{palettebullets}
\item \textbf{Lumi\`ere Lunaire} --- 月华铬黄是画面的「爆点」,月亮与星芒在蓝色重压下迸裂的光爆,是生命张力的极致
\item \textbf{Lueurs du Village} --- 村灯琥珀是画中唯一属于人间的暖,与冰冷的蓝色星空构成全画唯一的冷暖对峙
\item \textbf{Aube Glac\'ee} --- 晨曦苍蓝是过渡色,是深夜与月光交界处最微妙的呼吸感,介于压迫与解脱之间
\item \textbf{Tourbillon Outremer} --- 星涡群青是漩涡主体,是后印象派用色的标志——非自然色而是情绪色,是这组配色的「主旋」
\item \textbf{Cypr\`es Nocturne} --- 柏影墨绿是画面唯一「拒绝光」的存在,几近于黑,是视觉锚点,以最暗的色撑起整幅画的对比
\item \textbf{Minuit Cobalt} --- 深夜钴蓝是夜空最深处的情绪底色,是梵高精神世界的深渊,所有蓝色由它生长而来
\end{palettebullets}
\par\bigskip
% ---------------------------------------------------------------------------
\bookpalettecardhead{14\enspace A Thousand Li}
\palettesubsec{色彩哲学}
\noindent 《千里江山图》由北宋王希孟以矿物颜料层层积染而成——石青、石绿研磨后反复叠压,非一次上色,而是以「薄中见厚」的积染工艺令色泽沉入绢底,形成一种由内向外透光的矿物质感。六色之间非独立存在,而是共同构成一套山水呼吸的色彩生态。

\begin{fullwidth}\small\noindent
\begin{tabularx}{\linewidth}{@{}clllX@{}}
\toprule\rowcolor{booktableheadcolor}
\booktableheadcell{} & \booktableheadcell{中文名} & \booktableheadcell{原名} & \booktableheadcell{RGB} & \booktableheadcell{释义} \\
\midrule
\bookcardswatch{218,203,170} & 宋绢暖 & \foreignname{宋绢} \textperiodcentered\ S\`ong-ju\=an & 218, 203, 170 & 千年绢素氧化后的暖米底色,非颜料而是时间本身的色彩 \\
\bookcardswatch{110,165,195} & 烟渚色 & \foreignname{烟渚} \textperiodcentered\ Y\=an-zh\v{u} & 110, 165, 195 & 石青与铅白调和后的淡化层,用于远山轮廓与江面折光 \\
\bookcardswatch{183,118,45} & 砂岩赭 & \foreignname{砂赭} \textperiodcentered\ Sh\=a-zh\v{e} & 183, 118, 45 & 赭石(褐铁矿)与泥金混调,用于勾勒礁石肌理、屋舍廊桥与水波金线 \\
\bookcardswatch{96,148,110} & 岚霭绿 & \foreignname{岚霭绿} \textperiodcentered\ L\'an-\v{a}i-l\"u & 96, 148, 110 & 石绿渐次减量后的过渡色,表现中景植被与山腰云雾交融处 \\
\bookcardswatch{48,105,76} & 苍岭绿 & \foreignname{苍岭绿} \textperiodcentered\ C\=ang-l\v{i}ng-l\"u & 48, 105, 76 & 天然孔雀石的正色,层叠于山峦受光面,画面最具生命感的色调 \\
\bookcardswatch{35,78,112} & 霁山青 & \foreignname{霁山青} \textperiodcentered\ J\`i-sh\=an-q\=ing & 35, 78, 112 & 天然蓝铜矿研细后的深沉矿蓝,画中山体背阴面与水深处的主色 \\
\bottomrule
\end{tabularx}\end{fullwidth}

\palettesubsec{五大特征的色彩映射}
\begin{palettebullets}
\item \textbf{青绿主调} --- 霁山青与苍岭绿共同锚定全组色彩的主旋,一冷一暖,一阴一阳,相互制衡而非竞争
\item \textbf{层次渐变} --- 从霁山青$\to$烟渚色、苍岭绿$\to$岚霭绿,各形成一条明度递进链,模拟积染工艺的「薄中见厚」
\item \textbf{金赭点缀} --- 砂岩赭以最小的面积承担最大的温度调节职责,令青绿主调不至于流于清冷
\item \textbf{空灵雅致} --- 烟渚色与岚霭绿的介入,在浓重矿色之间打开呼吸的空间,令整组色不堵不闷
\item \textbf{古绢质感} --- 宋绢暖作为底色托起五色,如同千年绢素将所有浓烈都沉淀为一种悠远的克制
\end{palettebullets}
\par\bigskip
% ---------------------------------------------------------------------------
\bookpalettecardhead{15\enspace And Quiet Flows the Don}
\palettesubsec{色彩哲学}
\noindent 肖霍洛夫用二十年写完这部史诗,顿河哥萨克的命运从不在鲜亮之处——它埋在冻土里、锈在马刀上、凝在战壕边的黑血里、散在大雪漫过的苇荡中。这组配色拒绝一切「颜色」,只留下时间磨损之后的残余:不是蓝,是浊蓝;不是红,是暗血;不是黄,是枯草;不是灰,是铅灰。

\begin{fullwidth}\small\noindent
\begin{tabularx}{\linewidth}{@{}clllX@{}}
\toprule\rowcolor{booktableheadcolor}
\booktableheadcell{} & \booktableheadcell{中文名} & \booktableheadcell{原名} & \booktableheadcell{RGB} & \booktableheadcell{释义} \\
\midrule
\bookcardswatch{142,120,70} & 枯苇秋黄 & \foreignname{Полынь} \textperiodcentered\ Polyn & 142, 120, 70 & 顿河边被霜打过的苇草与苦艾,褪色的黄绿 \\
\bookcardswatch{110,108,104} & 草原铅灰 & \foreignname{Степной Пепел} \textperiodcentered\ Stepnoy Pepel & 110, 108, 104 & 冬日荒原的天与地混为一色,哥萨克男人眼神里的那种灰 \\
\bookcardswatch{118,86,52} & 黄土大地 & \foreignname{Донская Земля} \textperiodcentered\ Donskaya Zemlya & 118, 86, 52 & 顿河流域特有的栗钙土,带腐殖质气息的褐色 \\
\bookcardswatch{122,66,36} & 铁锈赭红 & \foreignname{Ржавое Железо} \textperiodcentered\ Rzhavoye Zhelezo & 122, 66, 36 & 哥萨克军刀、马镫、步枪枪托上的锈蚀,战争腐蚀金属的颜色 \\
\bookcardswatch{52,70,90} & 顿河浊蓝 & \foreignname{Мутный Дон} \textperiodcentered\ Mutnyy Don & 52, 70, 90 & 顿河从不清澈——裹挟泥沙与战死者的血,全书的时间底色 \\
\bookcardswatch{88,26,24} & 暗血战痕 & \foreignname{Запёкшаяся Кровь} \textperiodcentered\ Zapyokshayasya Krov & 88, 26, 24 & 凝固三日后渗入军衣的深褐红,血色命运在这里沉默 \\
\bottomrule
\end{tabularx}\end{fullwidth}

\palettesubsec{色彩结构的悲剧逻辑}
\begin{palettebullets}
\item \textbf{冷灰压顶} --- 草原铅灰与顿河浊蓝构成整组色彩的天空与水面,将所有暖色死死压在下方,如同历史对个体命运的碾压从未松手
\item \textbf{土褐承重} --- 黄土大地与枯苇秋黄是这片土地仅剩的「生」意,但它们同样去饱和、去鲜亮,美也是疲惫的美
\item \textbf{暗血锚底} --- 暗血战痕是全组最深的极色,不做主角,只在最沉的阴影处出现——如同《静静的顿河》里的死亡,从不戏剧化
\item \textbf{锈感贯穿} --- 铁锈赭红横跨土褐与暗血之间,是这组配色的「时间感」载体,它告诉观者:这里所有的一切,都已经过去很久了
\end{palettebullets}
\par\bigskip
% ---------------------------------------------------------------------------
\bookpalettecardhead{16\enspace Cyberpunk Edgerunners}
\palettesubsec{色彩哲学}
\noindent 《边缘行者》的视觉语言本质是过载美学——色彩不是装饰,是症状。大卫的精神崩溃伴随着荧光色彩的失控侵蚀,夜之城的霓虹不是繁荣的象征而是麻醉剂,Lucy 的蓝是这个世界里唯一真实的出口。这组配色的核心矛盾是:极暗的底色与极亮的荧光并置,产生几近生理不适的振动感。

\begin{fullwidth}\small\noindent
\begin{tabularx}{\linewidth}{@{}clllX@{}}
\toprule\rowcolor{booktableheadcolor}
\booktableheadcell{} & \booktableheadcell{中文名} & \booktableheadcell{原名} & \booktableheadcell{RGB} & \booktableheadcell{释义} \\
\midrule
\bookcardswatch{225,255,8} & 迷幻荧黄 & \foreignname{サイコ黄} \textperiodcentered\ Saiko-ki & 225, 255, 8 & 大卫夹克的灵魂色,赛博精神病发作时视觉过载之色 \\
\bookcardswatch{10,238,100} & 电路荧绿 & \foreignname{電脳緑} \textperiodcentered\ Denn\=o{}-ryoku & 10, 238, 100 & 网络骇入的代码流与改造体的皮下光纹 \\
\bookcardswatch{238,18,120} & 霓虹品红 & \foreignname{ネオン桃} \textperiodcentered\ Neon-momo & 238, 18, 120 & 夜之城最泛滥的霓虹,繁荣谎言的颜色 \\
\bookcardswatch{205,20,35} & 鲜血赤红 & \foreignname{エッジ赤} \textperiodcentered\ Edge-aka & 205, 20, 35 & 边缘行者的代价,街头枪战后水泥地上的颜色 \\
\bookcardswatch{15,32,98} & 露娜深蓝 & \foreignname{月の蒼} \textperiodcentered\ Tsuki-no-ao & 15, 32, 98 & Lucy 仰望月球的眼神,唯一不属于夜之城的颜色 \\
\bookcardswatch{14,8,28} & 夜都虚空 & \foreignname{夜都虚空} \textperiodcentered\ Yato Kok\=u{} & 14, 8, 28 & 夜之城街道底层的绝对黑暗,含深紫底调 \\
\bottomrule
\end{tabularx}\end{fullwidth}

\palettesubsec{色彩结构的叙事逻辑}
\begin{palettebullets}
\item \textbf{虚空压底} --- 夜都虚空作为绝对底色,饱和度趋近于零,是让其余五色得以「爆炸」的必要前提。没有黑暗,霓虹只是颜料
\item \textbf{冷暖裂变} --- 露娜深蓝与霓虹品红、迷幻荧黄构成最极端的冷暖对立:蓝是 Lucy 的月球梦,粉与黄是夜之城的谎言与疯狂
\item \textbf{荧光过载} --- 迷幻荧黄与电路荧绿是赛博精神病美学的视觉核心,两色并置时产生几近生理不适的振动感
\item \textbf{红色终章} --- 鲜血赤红在色板中不做主角,却是最诚实的颜色——像结局一样,无可回避
\end{palettebullets}
\par\bigskip
% ---------------------------------------------------------------------------
\bookpalettecardhead{17\enspace The Grand Budapest Hotel}
\palettesubsec{色彩哲学}
\noindent 韦斯·安德森的配色从不是写实的,而是经过精密校准的情感滤镜——《布达佩斯大饭店》的粉色不是粉色,是一整个正在消逝的旧欧洲贵族世界被压缩成的颜色;紫色不是紫色,是古斯塔夫先生永不失礼的最后尊严。每一种颜色都经过刻意的去真实化处理。

\begin{fullwidth}\small\noindent
\begin{tabularx}{\linewidth}{@{}clllX@{}}
\toprule\rowcolor{booktableheadcolor}
\booktableheadcell{} & \booktableheadcell{中文名} & \booktableheadcell{原名} & \booktableheadcell{RGB} & \booktableheadcell{释义} \\
\midrule
\bookcardswatch{247,235,215} & 香草奶油 & Cr\`eme Vanille & 247, 235, 215 & 饭店内壁、Mendl 蛋糕胚、信笺底色,童话世界的空气本身 \\
\bookcardswatch{168,178,196} & 阿尔卑斯雾蓝 & Brume Alpine & 168, 178, 196 & 积雪山景、冬日车厢窗外的低饱和冷调,情绪降温剂 \\
\bookcardswatch{237,148,158} & 美酒蔷薇 & Rose M\'endl & 237, 148, 158 & 饭店外墙与 Mendl 甜品盒的标志粉,掺了一丝尘埃的旧玫瑰色 \\
\bookcardswatch{176,136,60} & 古铜暗金 & Dor\'e Antique & 176, 136, 60 & 镀金镜框、门把手、制服肩章的年岁感,旧贵族的光泽 \\
\bookcardswatch{128,88,148} & 执行者紫 & Violet Concierge & 128, 88, 148 & 古斯塔夫先生制服的紫,兼具权威与荒诞 \\
\bookcardswatch{143,52,65} & 复古酒绛 & Bordeaux Vintage & 143, 52, 65 & 深色场景的戏剧性锚点——图书馆皮革、红酒、密谋与危险 \\
\bottomrule
\end{tabularx}\end{fullwidth}

\palettesubsec{五重美学法则的色彩映射}
\begin{palettebullets}
\item \textbf{马卡龙柔色} --- 蔷薇粉与执行者紫是主视觉担当,高饱和但经过刻意的明度压制,像马卡龙饼壳——色彩分明却绝不刺眼
\item \textbf{奶油暖底} --- 香草奶油以最大留白统摄全局,是韦斯·安德森式对称构图里的「负空间」,令其他五色不至于拥挤
\item \textbf{复古暗金} --- 古铜暗金专门负责「贵族感」的重量,它不参与童话叙事,只在边框与细节处发光
\item \textbf{冷调克制} --- 阿尔卑斯雾蓝是全组唯一的冷色,低饱和度让它成为调节器,在甜腻的暖色系中打开一扇透气的窗
\item \textbf{甜而不腻} --- 复古酒绛是整组配色最精妙的设计:一粒深沉的暗色锚点,像甜点里的一滴苦杏仁酒,让所有的轻盈都有了重量
\end{palettebullets}
\par\bigskip
% ---------------------------------------------------------------------------
\bookpalettecardhead{18\enspace Renaissance Florence}
\palettesubsec{色彩哲学}
\noindent 佛罗伦萨的色彩并非设计而来,而是从物质本身生长出来的——青金石研磨成颜料、铅白调入蛋黄成坦培拉底料、金箔贴于木质祭坛、陶土在托斯卡纳窑火中烧成穹顶砖瓦。这六种颜色皆有真实的物质出处:它们曾是颜料、矿石、建材、木料,是波提切利、乔托与多纳泰罗共同呼吸的色谱。

\begin{fullwidth}\small\noindent
\begin{tabularx}{\linewidth}{@{}clllX@{}}
\toprule\rowcolor{booktableheadcolor}
\booktableheadcell{} & \booktableheadcell{中文名} & \booktableheadcell{原名} & \booktableheadcell{RGB} & \booktableheadcell{释义} \\
\midrule
\bookcardswatch{238,228,208} & 翡冷翠象牙 & Avorio Fiorentino & 238, 228, 208 & 卡拉拉大理石截面与坦培拉画板的打底铅白,含蛋黄与亚麻油的暖意 \\
\bookcardswatch{190,148,52} & 祭坛暖金 & Oro dell'Altare & 190, 148, 52 & 圣母像背光与美第奇金饰的金箔色,五百年空气氧化后内敛发光的旧金 \\
\bookcardswatch{172,84,50} & 穹顶陶红 & Cotto Brunellesco & 172, 84, 50 & 布鲁内莱斯基穹顶砖瓦与佛罗伦萨屋檐的标志陶土红 \\
\bookcardswatch{54,80,60} & 柏影墨绿 & Verde Cipresso & 54, 80, 60 & 托斯卡纳丘陵柏树与铜绿颜料(Verdigris)的混调 \\
\bookcardswatch{46,70,118} & 青金矿蓝 & Oltremare di Lapislazzuli & 46, 70, 118 & 产自阿富汗矿山的青金石研磨所得,价比黄金,专用于圣母袍服 \\
\bookcardswatch{85,54,34} & 胡桃深褐 & Noce Toscano & 85, 54, 34 & 佛罗伦萨工坊核桃木画板与深色底釉的颜色 \\
\bottomrule
\end{tabularx}\end{fullwidth}

\palettesubsec{四重气质的色彩映射}
\begin{palettebullets}
\item \textbf{神圣庄重} --- 青金矿蓝承担全组的精神重量。青金石在文艺复兴时期价比黄金,非普通委托可用,它的存在本身即是神圣的声明;祭坛暖金以最克制的方式散发光晕,二者共同构建宗教艺术的仪式感
\item \textbf{温润华贵} --- 翡冷翠象牙与祭坛暖金是贵族底色的两面——一个是触感温润的大理石与绢布,一个是美第奇家族以财富供养艺术的金属记忆;陶土红以人间温度调和神圣的冷峻
\item \textbf{古典醇厚} --- 六色均经过刻意的去现代化处理:无一纯色,皆含矿物杂质的「不纯粹」——这正是天然颜料与合成色料最根本的区别。厚重来自饱和度的克制,醇厚来自每个颜色内部的复杂性
\item \textbf{人文雅致} --- 柏影墨绿与胡桃深褐是人文主义的色彩锚点:绿来自托斯卡纳的自然山丘,褐来自工坊的核桃木画板,它们将神圣从天空拉回大地
\end{palettebullets}
\par\bigskip
% ---------------------------------------------------------------------------
\bookpalettecardhead{19\enspace Soviet Avant-Garde}
\palettesubsec{色彩哲学}
\noindent 苏联先锋派从不「设计配色」——他们宣告颜色。罗钦科说色彩必须有任务,李西茨基说形式即政治,马列维奇说至上主义的黑色正方形是零度的绝对。这组配色拒绝调和、拒绝渐变、拒绝装饰性的过渡,每种颜色都是一个硬边宣言:红色不解释自己,黑色不道歉,灰色不讨好任何人。

\begin{fullwidth}\small\noindent
\begin{tabularx}{\linewidth}{@{}clllX@{}}
\toprule\rowcolor{booktableheadcolor}
\booktableheadcell{} & \booktableheadcell{中文名} & \booktableheadcell{原名} & \booktableheadcell{RGB} & \booktableheadcell{释义} \\
\midrule
\bookcardswatch{225,218,205} & 粗纸白 & \foreignname{ГАЗЕТА} \textperiodcentered\ Gazeta & 225, 218, 205 & 《真理报》与构成主义宣传册的新闻纸白,含油墨渗透后的暖灰底调 \\
\bookcardswatch{200,150,32} & 宣传赭黄 & \foreignname{ПЛАКАТ} \textperiodcentered\ Plakat & 200, 150, 32 & 克鲁齐斯摄影蒙太奇海报与农业集体化宣传画的赭黄 \\
\bookcardswatch{118,114,108} & 混凝土灰 & \foreignname{БЕТОН} \textperiodcentered\ Beton & 118, 114, 108 & 塔特林塔、莫斯科工厂厂房、五年计划基建工地的浇筑混凝土色 \\
\bookcardswatch{50,80,138} & 机械蓝 & \foreignname{ЧЕРТЁЖ} \textperiodcentered\ Chertyozh & 50, 80, 138 & 工程蓝图与苏联工人阶级制服的冷蓝,去饱和处理 \\
\bookcardswatch{196,28,28} & 列宁红 & \foreignname{КРАСНЫЙ} \textperiodcentered\ Krasnyy & 196, 28, 28 & 布尔什维克旗帜与罗钦科海报的革命红,刻意压低饱和度令其沉重如铁 \\
\bookcardswatch{24,20,18} & 铸铁黑 & \foreignname{ЧЁРНЫЙ} \textperiodcentered\ Chyornyy & 24, 20, 18 & 马列维奇《黑色正方形》的绝对黑,虚无与确定性同时存在的黑 \\
\bottomrule
\end{tabularx}\end{fullwidth}

\palettesubsec{六重美学法则的色彩映射}
\begin{palettebullets}
\item \textbf{红黑强撞} --- 列宁红与铸铁黑是整组配色的结构核心,二者之间无过渡、无缓冲,直接对撞——对比本身即是论点
\item \textbf{工业灰的中轴} --- 混凝土灰不参与任何戏剧,它只是站在那里承重,如同五年计划中那些无名工人,是这组配色的脊柱而非焦点
\item \textbf{机械蓝的精确} --- 机械蓝刻意压低至近乎去色的状态,与列宁红形成冷暖的意识形态对立:红是革命的激情,蓝是建设的纪律
\item \textbf{粗纸白的民主性} --- 新闻纸白是这组配色里唯一的「反精英」色,它拒绝象牙白的贵族温润,选择廉价印刷品的粗糙质感
\item \textbf{赭黄的未竟性} --- 宣传赭黄是全组配色中最具悲剧性的颜色:它是海报上的丰收麦穗、是乌托邦宣传画里永远灿烂的明天,是一个承诺的颜色而非现实的颜色
\end{palettebullets}
\par\bigskip
% ---------------------------------------------------------------------------
\bookpalettecardhead{20\enspace Constantinople}
\palettesubsec{色彩哲学}
\noindent 君士坦丁堡是人类历史上唯一一座同时属于两个大陆、三种文明、四种宗教的城市——它的颜色从不单纯。圣索菲亚大教堂的穹顶金光里叠压着希腊工匠、罗马工程、东方神学;博斯普鲁斯海峡的深蓝里混着黑海的寒流与爱琴海的温暖;帝王紫是用地中海骨螺提炼出的奢华。

\begin{fullwidth}\small\noindent
\begin{tabularx}{\linewidth}{@{}clllX@{}}
\toprule\rowcolor{booktableheadcolor}
\booktableheadcell{} & \booktableheadcell{中文名} & \booktableheadcell{原名} & \booktableheadcell{RGB} & \booktableheadcell{释义} \\
\midrule
\bookcardswatch{236,225,207} & 普罗科尼索象牙 & \foreignname{Λευκός} \textperiodcentered\ Fildi\c{s}i & 236, 225, 207 & 马尔马拉海小岛出产的普罗科尼索大理石,铺满圣索菲亚内殿地面 \\
\bookcardswatch{195,150,42} & 圣索菲亚金 & \foreignname{Χρυσός} \textperiodcentered\ Alt\i n & 195, 150, 42 & 圣索菲亚穹顶马赛克金箔的光,一千五百年烛烟熏染后内敛发光的旧金 \\
\bookcardswatch{170,112,50} & 狄奥多西赭 & \foreignname{\'{Ω}χρα} \textperiodcentered\ Toprak & 170, 112, 50 & 狄奥多西城墙砖石与博斯普鲁斯岸边老建筑的赭土色 \\
\bookcardswatch{50,76,58} & 柏树暗绿 & \foreignname{Κυπαρίσσι} \textperiodcentered\ Selvi & 50, 76, 58 & 君士坦丁堡山丘柏树与拜占庭绿色斑岩柱础的深沉暗绿 \\
\bookcardswatch{102,32,65} & 帝王骨螺紫 & \foreignname{Πορφύρα} \textperiodcentered\ Mor & 102, 32, 65 & 拜占庭皇帝专属的骨螺紫,提炼一克需杀万只骨螺 \\
\bookcardswatch{30,52,88} & 博斯普鲁斯蓝 & \foreignname{Βόσπορος} \textperiodcentered\ Bo\u{g}az & 30, 52, 88 & 海峡傍晚收拢光线后的深沉蓝,黑海与爱琴海在此交汇 \\
\bottomrule
\end{tabularx}\end{fullwidth}

\palettesubsec{文明交汇的色彩结构}
\begin{palettebullets}
\item \textbf{东西冷暖轴} --- 博斯普鲁斯深蓝(西方理性、海洋秩序)与狄奥多西赭(东方大地、安纳托利亚土壤)构成配色的东西文明轴线;二者温度相反,却在君士坦丁堡这座城市里共存了千年而未曾消解
\item \textbf{神圣三角} --- 圣索菲亚金、帝王骨螺紫、普罗科尼索象牙共同构建宗教神圣感的色彩空间:金是神的光,紫是神在人间的代理人,象牙是神居住其中的建筑肌肤
\item \textbf{紫红的文明桥接} --- 帝王骨螺紫在整组配色中承担最复杂的角色:它同时属于罗马帝国的皇权传统、拜占庭基督教的神学色彩,以及波斯-东方宫廷的奢华美学
\item \textbf{暗绿的历史纵深} --- 柏树暗绿以最低调的姿态承载最漫长的时间跨度:它同时出现在拜占庭墓地柏树、奥斯曼庭园绿化与伊斯兰清真寺的铜绿穹顶上
\item \textbf{金辉内敛} --- 整组配色刻意将圣索菲亚金的饱和度压低至矿物感区间,令其发光而不炫目——这是拜占庭马赛克艺术的核心逻辑:金不是装饰,是光本身的物质化
\end{palettebullets}
\par\bigskip
% ---------------------------------------------------------------------------
\bookpalettecardhead{21\enspace France}
\palettesubsec{色彩哲学}
\noindent 法国的颜色从不是单纯的颜色——它是一种世界观的物质显现。共和红里藏着断头台与玫瑰,深蓝里住着笛卡尔的理性与莫奈的晨雾,香槟白是凡尔赛大理石与左岸咖啡馆信纸同一种底色。这组配色试图捕捉法国精神最核心的内在矛盾:极度理性与极度浪漫在同一块土地上共生。

\begin{fullwidth}\small\noindent
\begin{tabularx}{\linewidth}{@{}clllX@{}}
\toprule\rowcolor{booktableheadcolor}
\booktableheadcell{} & \booktableheadcell{中文名} & \booktableheadcell{原名} & \booktableheadcell{RGB} & \booktableheadcell{释义} \\
\midrule
\bookcardswatch{180,32,40} & 共和红 & Rouge Marianne & 180, 32, 40 & 三色旗之红,但非正红——波尔多红酒的深度与自由女神帽的激情的中间值 \\
\bookcardswatch{232,222,205} & 香槟象牙 & Ivoire Champagne & 232, 222, 205 & 非纯白——香槟气泡的暖调、卢浮宫石灰岩的温度、左岸咖啡馆信纸 \\
\bookcardswatch{44,74,138} & 共和深蓝 & Bleu R\'epublique & 44, 74, 138 & 三色旗之蓝,取塞夫勒瓷器皇家蓝的深度——启蒙理性的蓝 \\
\bookcardswatch{143,108,148} & 普罗旺斯薰衣草 & Lavande Proven\c{c}ale & 143, 108, 148 & 南法薰衣草田在午后阳光与灰白石灰土之间呈现的真实颜色 \\
\bookcardswatch{142,133,58} & 凡尔赛橄榄金 & Olive Dor\'ee & 142, 133, 58 & 凡尔赛宫正式花园修剪黄杨篱的金绿——古典法式园林美学 \\
\bookcardswatch{122,126,132} & 巴黎锌灰 & Zinc Parisien & 122, 126, 132 & 奥斯曼改造后巴黎屋顶锌板在阴天的反光色 \\
\bottomrule
\end{tabularx}\end{fullwidth}

\palettesubsec{法兰西精神的色彩结构}
\begin{palettebullets}
\item \textbf{三色旗的再解读} --- 共和红、香槟象牙、共和深蓝并非原版三色旗的复刻,而是对三色精神的去符号化再提炼:红被赋予波尔多的重量,白被赋予时间的温度,蓝被赋予启蒙的深度——旗帜是政治的,这组色是文化的
\item \textbf{南北气质的张力} --- 普罗旺斯薰衣草与巴黎锌灰构成法国南北最核心的精神对立:南法是感官、土地、阳光与慵懒的浪漫;巴黎是理性、几何、灰色天空与克制的优雅
\item \textbf{先锋撞色的潜伏} --- 凡尔赛橄榄金与共和深蓝的并置是全组最具先锋张力的组合:一个属于太阳王的古典秩序,一个属于共和革命的理性锋芒——色相接近却语义对立
\item \textbf{克制高级的实现路径} --- 六色均经过刻意的饱和度压制:红不燃烧,蓝不冰冷,紫不妖艳,金不张扬,灰不死板,白不刺目。法国人所理解的「高级」,从来不是色彩的强度,而是色彩恰到好处地收敛自身
\item \textbf{土地温暖的隐线} --- 香槟象牙与普罗旺斯薰衣草共同编织这组配色的「土地感」底线:前者是石灰岩与小麦的暖,后者是干燥薰衣草与灰白土壤的暖——法国的浪漫从不悬浮在空中
\end{palettebullets}
\par\bigskip
% ---------------------------------------------------------------------------
\bookpalettecardhead{22\enspace Kyoto}
\palettesubsec{色彩哲学}
\noindent 京都的颜色从不主动开口——它们等待被看见。禅院的苔绿不争,枯山水的白砂不语,町家杉木的茶褐在雨中沉默,暗绀的夜空压着千年古都不作声。唯有朱红鸟居,在一切克制之中以极小的面积燃烧最大的温度,像茶室里一枝单花,冷寂里的那一点热烈。

\begin{fullwidth}\small\noindent
\begin{tabularx}{\linewidth}{@{}clllX@{}}
\toprule\rowcolor{booktableheadcolor}
\booktableheadcell{} & \booktableheadcell{中文名} & \booktableheadcell{原名} & \booktableheadcell{RGB} & \booktableheadcell{释义} \\
\midrule
\bookcardswatch{234,226,212} & 绢白 & \foreignname{絹白} \textperiodcentered\ Kinushiro & 234, 226, 212 & 西阵织绢布与和纸漉纸的底色,蚕丝与楮树皮的天然暖调 \\
\bookcardswatch{110,84,64} & 枯茶褐 & \foreignname{枯茶} \textperiodcentered\ Karacha & 110, 84, 64 & 町家百年杉木柱在时间与炉烟中形成的褐色,千利休「侘び」的颜色 \\
\bookcardswatch{204,184,178} & 樱灰粉 & \foreignname{桜鼠} \textperiodcentered\ Sakura-nezumi & 204, 184, 178 & 落樱覆于石灯笼上被雨水晕开后的颜色,物哀的精确定义 \\
\bookcardswatch{116,126,120} & 青竹灰 & \foreignname{青竹鼠} \textperiodcentered\ Aotake-nezumi & 116, 126, 120 & 岚山竹林晨雾中的竹节色,冷暖平衡的禅理体现 \\
\bookcardswatch{80,94,68} & 苔色 & \foreignname{苔色} \textperiodcentered\ Koke-iro & 80, 94, 68 & 龙安寺枯山水石组底部苔藓的深绿,时间留下的暗哑印记 \\
\bookcardswatch{34,40,66} & 紺 & \foreignname{紺} \textperiodcentered\ Kon & 34, 40, 66 & 能剧服装与京友禅夜色图案的深靛蓝,冥想状态本身的颜色 \\
\bottomrule
\end{tabularx}\end{fullwidth}

\palettesubsec{京都精神的色彩结构}
\begin{palettebullets}
\item \textbf{留白为骨} --- 绢白与樱灰粉共同构建配色的呼吸空间,前者是空无、负形、茶室墙面,后者是将消未消的感知边界——京都美学不在于填满,而在于知道哪里应当什么都没有
\item \textbf{冷暖天平} --- 六色底中,枯茶褐与苔绿属暖调,暗绀与青竹灰属冷调,樱灰粉与绢白居中游移——四季平衡,禅院冷寉与町家温暖在此精确对位
\item \textbf{幽玄的低饱和逻辑} --- 所有色彩均经过刻意的去鲜亮处理,模拟京都特有的盆地晨雾对色彩的柔化效果——雾不消色,只让色彩向内收,这正是「幽玄」的物理机制
\item \textbf{水墨留白的东方结构} --- 整组配色可理解为一幅水墨:暗绀是浓墨,苔绿是中墨,枯茶褐是淡墨,青竹灰是宿墨,樱灰粉是墨的晕散,绢白是纸——朱红是画师最后蘸朱砂落下的那一点钤印
\end{palettebullets}
\par\bigskip
% ---------------------------------------------------------------------------
\bookpalettecardhead{23\enspace Siamese Dream}
\palettesubsec{色彩哲学}
\noindent 《Siamese Dream》的封面本身就是一张配色宣言——两个女孩在柔焦的暖白光晕里相依,背景的深绿若隐若现,橙黄的光线像记忆里某个下午,美好得不真实。这不是摇滚专辑封面惯用的高对比冲击美学,而是一种刻意制造的梦境质感:过曝、柔化、带着胶片颗粒的温柔失真——如同 Corgan 把最原始的情感录进了最失真的吉他音墙里。

\begin{fullwidth}\small\noindent
\begin{tabularx}{\linewidth}{@{}clllX@{}}
\toprule\rowcolor{booktableheadcolor}
\booktableheadcell{} & \booktableheadcell{中文名} & \booktableheadcell{原名} & \booktableheadcell{RGB} & \booktableheadcell{释义} \\
\midrule
\bookcardswatch{238,226,208} & 终章暖白 & Luna & 238, 226, 208 & 封面的暖白光晕,专辑最后一首歌的余韵——不是纯白,是胶片过曝后仍保留的一丝人间温度,是梦醒前最后的柔光 \\
\bookcardswatch{182,180,193} & 裸露灰紫 & Disarm & 182, 180, 193 & 洗白的灰紫,Disarm 是整张专辑里防御最低的时刻,弦乐与铃声剥去了所有失真——这是玻璃碎裂前的颜色,透明、脆、带着童年创伤的寒意 \\
\bookcardswatch{198,158,65} & 可疑旧金 & Today & 198, 158, 65 & 沾了尘的旧金色,Today 听起来是全碟最明媚的旋律,却是 Corgan 最黑暗时刻的产物——这抹金因此必须是褪色的、带着微苦的甜,美丽而不可信任 \\
\bookcardswatch{150,145,135} & 悬浮烟灰 & Hummer & 150, 145, 135 & 烟灰褐,Hummer 用层层叠叠的音墙制造一种失重的悬浮感——这是梦游者眼中的颜色,既不清醒也不沉睡,是 shoegaze 美学最诚实的物质形态 \\
\bookcardswatch{188,105,40} & 低温烧橙 & Mayonaise & 188, 105, 40 & 深而沉的烧橙,封面光源的颜色——Mayonaise 是全碟最慢、最痛、最私密的曲目,这抹橙不是热情而是积压已久的情感在低温中缓慢焦化的颜色 \\
\bookcardswatch{42,68,48} & 深渊暗绿 & Soma & 42, 68, 48 & 封面背景若隐若现的深绿,九分钟的史诗长曲从轻柔攀爬至崩塌——这是有机物在黑暗中生长的颜色,苔藓与腐叶,美丽与窒息共生 \\
\bottomrule
\end{tabularx}\end{fullwidth}

\palettesubsec{青春期情感的悖论结构}
\begin{palettebullets}
\item \textbf{梦境包裹与失真美学} --- Luna 的暖白与 Hummer 的烟灰共同构成专辑的“失焦层”——一切痛苦都先被过曝、被柔化、被音墙包住,然后才允许被听见。这不是回避痛苦,而是让痛苦以可承受的密度释放
\item \textbf{美丽容器装最重的内容} --- Today 的旧金与 Disarm 的灰紫是悖论的两极:前者用明媚旋律包裹自毁倾向,后者用童谣式的铃声讲述童年创伤——Corgan 始终用最美的容器装最重的内容,这是 Siamese Dream 最核心的情感策略
\item \textbf{情感核心的低温焦化} --- Mayonaise 的烧橙是整组配色的情感核心(pal5=primary 也正落在此处)——不是燃烧的红,而是已经燃烧过、被时间慢烤、被自身重量压低的橙。它告诉读者:这里的热度不来自爆发,而来自长期不灭
\item \textbf{深渊作为最深锚点} --- Soma 的暗绿(pal6=accent)是六色中唯一的“黑色替身”——它不是纯黑,而是有机物的黑,是苔藓、腐叶、深林。整组配色因此从不依赖纯黑来收尾,而是用一种“仍在生长的黑暗”把整张专辑稳稳压住
\item \textbf{回避高饱和与强对比的渗透美学} --- 六色全部退避鲜亮,因为 Siamese Dream 的情绪从不“爆炸”——它“渗透”。没有任何一色试图夺取注意力,反而每一色都在以最低的姿态长时间地停留在视觉里,如同 shoegaze 的吉他延音
\end{palettebullets}


\end{document}